\worldchapter{\trans{weapons_title}}{appendix-weapons}
\label{ch:appendix-weapons}

\begin{multicols}{2}


\section{Klingenwaffen}



\subsection*{Stockdegen}

\begin{pexample}
Dieses Objekt erscheint wie ein modischer oder schlichter Gehstock, typischerweise aus poliertem Holz mit einer metallbeschlagenen Spitze und einem Ziergriff. Der Griff, oft ein Metallknauf oder L-förmiger Haken, lässt sich vom Hauptschaft abziehen. Dieser Schaft ist hohl und dient als Scheide für eine schlanke, steife und nadelspitze Stoßklinge, die im Inneren verborgen ist.
\end{pexample}

    {\small Nahkampf}\\

\begin{tabular}{@{}ll@{}}
    \textbf{\trans{weight_label}:} & 2.00 kg \\
    \textbf{\trans{price_label}:} & 160 \\
    \textbf{Durchschlag} & 1\\
    \textbf{Schadenspotential} & 1\\
\end{tabular}

\wbif{phasesix}{
    \vspace{3mm}
\faIcon{chess-rook}\hspace{0.5em}\faIcon{umbrella}\hspace{0.5em}\faIcon{flag}\hspace{0.5em}\faIcon{plane}\hspace{0.5em}\faIcon{microchip}\hspace{0.5em}}{}



\subsection*{Giftdolch}

\begin{pexample}
Dieser Dolch besitzt eine starre Klinge mit einer Einkerbung zum Auftragen von Gift, die in einer Spitze endet, welche besonders gut Rüstungen durchschlagen kann. Die Waffe ist dafür bekannt, ein Ziel in der Stärke des verwendeten Giftes zu vergiften.
\end{pexample}

    {\small Nahkampf}\\

\begin{tabular}{@{}ll@{}}
    \textbf{\trans{weight_label}:} & 0.80 kg \\
    \textbf{\trans{price_label}:} & 1400 \\
    \textbf{Durchschlag} & 2\\
    \textbf{Verborgenheit} & 2\\
    \textbf{Vorbereitung} & 0\\
    \textbf{Giftkanal} & 1\\
\end{tabular}

\wbif{phasesix}{
    \vspace{3mm}
\faIcon{eye}\hspace{0.5em}\faIcon{chess-rook}\hspace{0.5em}\faIcon{umbrella}\hspace{0.5em}\faIcon{flag}\hspace{0.5em}\faIcon{plane}\hspace{0.5em}\faIcon{microchip}\hspace{0.5em}}{}



\subsection*{Kurzschwert}

\begin{pexample}
Eine gerade, zweischneidige Klinge, sichtlich kürzer als ein Langschwert, aber deutlich länger als ein Dolch, bildet diese Waffe. Sie ist mit einem einfachen, einhändigen Griff, einer funktionalen Parierstange und einem simplen Knauf ausgestattet. Die Waffe ist leicht und ausbalanciert und dient eindeutig als praktische militärische Seitenwaffe.
\end{pexample}

    {\small Nahkampf}\\

\begin{tabular}{@{}ll@{}}
    \textbf{\trans{weight_label}:} & 0.70 kg \\
    \textbf{\trans{price_label}:} & 1000 \\
    \textbf{Durchschlag} & 1\\
    \textbf{Verborgenheit} & 2\\
    \textbf{Schadenspotential} & 1\\
\end{tabular}

\wbif{phasesix}{
    \vspace{3mm}
\faIcon{chess-rook}\hspace{0.5em}}{}



\subsection*{Anderthalbhänder}

\begin{pexample}
Diese gerade, zweischneidige Klinge ist merklich länger als ein gewöhnliches Ritterschwert, jedoch kürzer als ein echter Zweihänder. Ihr entscheidendes Merkmal ist der verlängerte Griff, der lang genug ist, um eine zweite Hand für erhöhte Kraft aufzunehmen, obwohl die Waffe für den einhändigen Gebrauch ausbalanciert bleibt.
\end{pexample}

    {\small Nahkampf}\\

\begin{tabular}{@{}ll@{}}
    \textbf{\trans{weight_label}:} & 1.20 kg \\
    \textbf{\trans{price_label}:} & 850 \\
    \textbf{Durchschlag} & 1\\
    \textbf{Verborgenheit} & 3\\
    \textbf{Schadenspotential} & 2\\
\end{tabular}

\wbif{phasesix}{
    \vspace{3mm}
\faIcon{chess-rook}\hspace{0.5em}\faIcon{umbrella}\hspace{0.5em}}{}



\subsection*{Florett}

\begin{pexample}
Eine außergewöhnlich dünne und flexible Klinge, oft mit viereckigem Querschnitt, definiert diese leichte Waffe. Sie besitzt keine geschärften Schneiden und läuft stattdessen in einer feinen, schmalen Spitze aus. Die Hand ist durch eine markante, schalenförmige Metallglocke geschützt.
\end{pexample}

    {\small Nahkampf}\\

\begin{tabular}{@{}ll@{}}
    \textbf{\trans{weight_label}:} & 1.30 kg \\
    \textbf{\trans{price_label}:} & 180 \\
    \textbf{Durchschlag} & 1\\
    \textbf{Verborgenheit} & 5\\
    \textbf{Schadenspotential} & 2\\
\end{tabular}

\wbif{phasesix}{
    \vspace{3mm}
\faIcon{chess-rook}\hspace{0.5em}\faIcon{umbrella}\hspace{0.5em}\faIcon{flag}\hspace{0.5em}\faIcon{plane}\hspace{0.5em}\faIcon{microchip}\hspace{0.5em}}{}



\subsection*{Schnappmesser}

\begin{pexample}
Ein Metallknopf oder Schalter ist direkt in die Seite des Griffs dieses Klappmessers eingelassen. Das Betätigen dieses Mechanismus löst eine interne Feder aus, wodurch die verborgene Klinge automatisch aufspringt, indem sie entweder seitlich herausklappt oder gerade aus der Oberseite hervorschießt.
\end{pexample}

    {\small Nahkampf}\\

\begin{tabular}{@{}ll@{}}
    \textbf{\trans{weight_label}:} & 0.30 kg \\
    \textbf{\trans{price_label}:} & 400 \\
    \textbf{Verborgenheit} & 1\\
    \textbf{Vorbereitung} & 0\\
    \textbf{Schadenspotential} & 2\\
\end{tabular}

\wbif{phasesix}{
    \vspace{3mm}
\faIcon{flag}\hspace{0.5em}\faIcon{plane}\hspace{0.5em}\faIcon{microchip}\hspace{0.5em}}{}



\subsection*{Zweihänder}

\begin{pexample}
Eine enorme, gerade, zweischneidige Klinge, weit länger als ein gewöhnliches Ritterschwert, dominiert diese Waffe. Ihr entscheidendes Merkmal ist das massive, zweihändige Heft, das einen verlängerten Griff und einen schweren Knauf als Gegengewicht nutzt. Die Waffe ist mit einer breiten, funktionalen Parierstange ausgestattet, die auf ihre immense Größe ausgelegt ist.
\end{pexample}

    {\small Nahkampf}\\

\begin{tabular}{@{}ll@{}}
    \textbf{\trans{weight_label}:} & 2.50 kg \\
    \textbf{\trans{price_label}:} & 4000 \\
    \textbf{Durchschlag} & 1\\
    \textbf{Reichweite} & 2\\
    \textbf{Verborgenheit} & 8\\
    \textbf{Schadenspotential} & 3\\
\end{tabular}

\wbif{phasesix}{
    \vspace{3mm}
\faIcon{cube}\hspace{0.5em}}{}



\subsection*{Katana}

\begin{pexample}
Eine lange, schlanke Klinge mit einer ausgeprägten, einschneidigen Krümmung definiert dieses Schwert. Es ist mit einem für den zweihändigen Griff ausgelegten Griff montiert, der oft mit Kordel über Rochenhaut gewickelt ist. Ein markantes rundes oder eckiges Metallstichblatt (Tsuba) trennt den Griff von der Klinge.
\end{pexample}

    {\small Nahkampf}\\

\begin{tabular}{@{}ll@{}}
    \textbf{\trans{weight_label}:} & 0.80 kg \\
    \textbf{\trans{price_label}:} & 1200 \\
    \textbf{Durchschlag} & 1\\
    \textbf{Verborgenheit} & 3\\
    \textbf{Schadenspotential} & 3\\
\end{tabular}

\wbif{phasesix}{
    \vspace{3mm}
\faIcon{cube}\hspace{0.5em}}{}



\subsection*{Katar}

\begin{pexample}
Der Katar, der in abliegenden Regionen bekannt ist, wird durch ihren einzigartigen H-förmigen Metallgriff definiert. Der Griff besteht aus zwei parallelen Stangen, welche die Seiten der Hand schützen und durch einen oder zwei Querbalken für die Faust verbunden sind. Eine breite, spitze, zweischneidige Klinge ist an der Vorderseite dieses Griffs befestigt und ragt geradewegs von den Knöcheln des Trägers nach vorne.
\end{pexample}

    {\small Nahkampf}\\

\begin{tabular}{@{}ll@{}}
    \textbf{\trans{weight_label}:} & 0.80 kg \\
    \textbf{\trans{price_label}:} & 450 \\
    \textbf{Durchschlag} & 1\\
    \textbf{Verborgenheit} & 2\\
    \textbf{Schadenspotential} & 2\\
\end{tabular}

\wbif{phasesix}{
    \vspace{3mm}
\faIcon{chess-rook}\hspace{0.5em}\faIcon{umbrella}\hspace{0.5em}\faIcon{flag}\hspace{0.5em}\faIcon{plane}\hspace{0.5em}\faIcon{microchip}\hspace{0.5em}}{}



\subsection*{Klingenbrecher}

\begin{pexample}
Wird mit dieser Waffe erfolgreich angegriffen oder pariert, nimmt die Waffe des Gegners für jeden Erfolg 10\% Schaden, zusätzlich zu den Wunden.
\end{pexample}

    {\small Nahkampf}\\

\begin{tabular}{@{}ll@{}}
    \textbf{\trans{weight_label}:} & 1.20 kg \\
    \textbf{\trans{price_label}:} & 850 \\
    \textbf{Verborgenheit} & 3\\
    \textbf{Schadenspotential} & 1\\
\end{tabular}

\wbif{phasesix}{
    \vspace{3mm}
\faIcon{chess-rook}\hspace{0.5em}\faIcon{umbrella}\hspace{0.5em}\faIcon{flag}\hspace{0.5em}\faIcon{plane}\hspace{0.5em}\faIcon{microchip}\hspace{0.5em}}{}



\subsection*{Entermesser}

\begin{pexample}
Oft als Waffe eines Piraten bezeichnet, verfügt dieses Schwert über eine kurze, breite und leicht gekrümmte einschneidige Klinge. Der Griff wird von einem massiven, oft schmucklosen, korb- oder schalenförmigen Stichblatt geschützt, das die Hand umschließt. Der Stahl wirkt dick und funktional, ihm fehlt der feine Schliff eines hochwertigen Schwertes, was auf eine geringerwertige, zweckdienliche Konstruktion hindeutet.
\end{pexample}

    {\small Nahkampf}\\

\begin{tabular}{@{}ll@{}}
    \textbf{\trans{weight_label}:} & 1.20 kg \\
    \textbf{\trans{price_label}:} & 1000 \\
    \textbf{Verborgenheit} & 4\\
    \textbf{Schadenspotential} & 3\\
\end{tabular}

\wbif{phasesix}{
    \vspace{3mm}
\faIcon{eye}\hspace{0.5em}\faIcon{chess-rook}\hspace{0.5em}\faIcon{umbrella}\hspace{0.5em}\faIcon{flag}\hspace{0.5em}\faIcon{plane}\hspace{0.5em}\faIcon{microchip}\hspace{0.5em}}{}



\subsection*{Langes Messer}

\begin{pexample}
Die Klinge des langen Messers ist lang, gerade und einschneidig, merklich größer als ein Standarddolch, besitzt jedoch nicht die Parierstange und Balance eines echten Kurzschwertes. Der Griff ist einfach, oft nur aus Holz oder mit Leder umwickelt, und hat einen minimalen oder fehlenden Handschutz. Es ist eine praktische, beinahe grobe Waffe, die irgendwo zwischen einem großen Werkzeug und einer reinen Kampfklinge angesiedelt ist.
\end{pexample}

    {\small Nahkampf}\\

\begin{tabular}{@{}ll@{}}
    \textbf{\trans{weight_label}:} & 1.20 kg \\
    \textbf{\trans{price_label}:} & 300 \\
    \textbf{Verborgenheit} & 4\\
    \textbf{Schadenspotential} & 2\\
\end{tabular}

\wbif{phasesix}{
    \vspace{3mm}
\faIcon{cube}\hspace{0.5em}}{}



\subsection*{Skalpell}

\begin{pexample}
Das Skalpell besteht aus einem schlanken Griff aus Metall oder Knochen, der für einen feinen, präzisen Griff bemessen ist. Er hält eine kleine, feststehende Klinge von außergewöhnlicher Schärfe. Die Klinge selbst ist dünn und besitzt oft eine ausgeprägte Krümmung oder eine feine Spitze, wodurch sie weitaus zerbrechlicher und schneidender wirkt als ein Allzweckmesser.
\end{pexample}

    {\small Nahkampf}\\

\begin{tabular}{@{}ll@{}}
    \textbf{\trans{weight_label}:} & 0.50 kg \\
    \textbf{\trans{price_label}:} & 250 \\
    \textbf{Verborgenheit} & 1\\
    \textbf{Vorbereitung} & 0\\
    \textbf{Schadenspotential} & 2\\
\end{tabular}

\wbif{phasesix}{
    \vspace{3mm}
\faIcon{chess-rook}\hspace{0.5em}\faIcon{umbrella}\hspace{0.5em}\faIcon{flag}\hspace{0.5em}\faIcon{plane}\hspace{0.5em}\faIcon{microchip}\hspace{0.5em}}{}



\subsection*{Säbel}

\begin{pexample}
Eine lange, gekrümmte, einschneidige Klinge definiert dieses einhändige Schwert. Sein markantestes Merkmal ist der Handschutz, das über einen hervorstehenden Metall-Handbügel verfügt, der sich von der Parierstange bis zum Knauf erstreckt und die Hand umschließt und schützt. Die Klinge selbst ist oft mäßig breit, was sie vom schlanken Profil eines Rapiers unterscheidet.
\end{pexample}

    {\small Nahkampf}\\

\begin{tabular}{@{}ll@{}}
    \textbf{\trans{weight_label}:} & 1.00 kg \\
    \textbf{\trans{price_label}:} & 950 \\
    \textbf{Verborgenheit} & 3\\
    \textbf{Schadenspotential} & 3\\
\end{tabular}

\wbif{phasesix}{
    \vspace{3mm}
\faIcon{cube}\hspace{0.5em}}{}



\subsection*{Bajonett}

\begin{pexample}
Diese moderne Waffe besteht aus einer geraden, scharfen Stahlklinge, die einem langen Dolch ähnelt. Ihr Griff ist nicht zum Halten gedacht, sondern weist einen markanten hohlen Ring sowie einen Arretierungsmechanismus auf. Diese anbringbare Konstruktion, eindeutig für den militärischen Gebrauch bestimmt, ermöglicht es, den gesamten Gegenstand an den Lauf einer Langwaffe zu montieren.
\end{pexample}

    {\small Nahkampf}\\

\begin{tabular}{@{}ll@{}}
    \textbf{\trans{weight_label}:} & 1.80 kg \\
    \textbf{\trans{price_label}:} & 850 \\
    \textbf{Durchschlag} & 1\\
    \textbf{Reichweite} & 2\\
    \textbf{Verborgenheit} & 4\\
    \textbf{Schadenspotential} & 2\\
\end{tabular}

\wbif{phasesix}{
    \vspace{3mm}
\faIcon{umbrella}\hspace{0.5em}\faIcon{flag}\hspace{0.5em}}{}



\subsection*{Molekularmesser}

\begin{pexample}
Diese Waffe ähnelt einem Kampfmessergriff aus mattschwarzen Polymeren, doch anstelle einer konventionellen Klinge besitzt sie einen kurzen, starren Grat. Dieser Grat stützt eine schimmernde, monomolekulare Schneide, die unmöglich dünn und vollkommen gleichmäßig erscheint und keinerlei sichtbare Zahnung oder Schliff aufweist.
\end{pexample}

    {\small Nahkampf}\\

\begin{tabular}{@{}ll@{}}
    \textbf{\trans{weight_label}:} & 0.30 kg \\
    \textbf{\trans{price_label}:} & 1000 \\
    \textbf{Durchschlag} & 2\\
    \textbf{Verborgenheit} & 1\\
    \textbf{Vorbereitung} & 0\\
    \textbf{Schadenspotential} & 2\\
\end{tabular}

\wbif{phasesix}{
    \vspace{3mm}
\faIcon{space-shuttle}\hspace{0.5em}}{}



\subsection*{Spatha}

\begin{pexample}
Die Spatha (Plural Spathae, lateinisch für „Breite Klinge“) ist ein zweischneidiges, vorwiegend zum Hieb konzipiertes, einhändig geführtes Schwert mit gerader Klinge. Diese Schwertform existierte etwa vom 1. Jahrhundert v. Chr. bis zum Ende des Frühmittelalters. Über die Jahrhunderte trat sie in verschiedenen Formen auf, insbesondere germanische Ringschwerter der Völkerwanderungszeit bis hin zu den volkstümlich so genannten „Wikingerschwertern“.
\end{pexample}

    {\small Nahkampf}\\

\begin{tabular}{@{}ll@{}}
    \textbf{\trans{weight_label}:} & 1.50 kg \\
    \textbf{\trans{price_label}:} & 250 \\
    \textbf{Durchschlag} & 1\\
    \textbf{Reichweite} & 2\\
    \textbf{Schadenspotential} & 2\\
\end{tabular}

\wbif{phasesix}{
    \vspace{3mm}
\faIcon{eye}\hspace{0.5em}}{}



\subsection*{Krummsäbel}

\begin{pexample}
Dieses einhändige Schwert wird durch seine ausgeprägte, schwungvolle Krümmung definiert und besitzt eine einzelne scharfe Schneide an der Außenseite. Die Klinge verbreitert sich oft leicht zur Spitze hin. Es ist mit einem einfachen Griff und einer Parierstange ausgestattet, die meist gerade oder leicht zur Klinge hin gebogen ist.
\end{pexample}

    {\small Nahkampf}\\

\begin{tabular}{@{}ll@{}}
    \textbf{\trans{weight_label}:} & 1.00 kg \\
    \textbf{\trans{price_label}:} & 500 \\
    \textbf{Verborgenheit} & 5\\
    \textbf{Schadenspotential} & 3\\
\end{tabular}

\wbif{phasesix}{
    \vspace{3mm}
\faIcon{chess-rook}\hspace{0.5em}\faIcon{umbrella}\hspace{0.5em}\faIcon{flag}\hspace{0.5em}\faIcon{plane}\hspace{0.5em}\faIcon{microchip}\hspace{0.5em}}{}



\subsection*{Lichtschwert}

\begin{pexample}
Ein zylindrischer Griff aus poliertem Metall, für ein oder zwei Hände bemessen, beherbergt die Komponenten dieser Waffe. Er verfügt über einen eingelassenen Aktivierungsknopf und eine markante Emitter-Öffnung an einem Ende. Im aktiven Zustand fährt ein solider, klingenlanger Strahl aus intensiver, glühender Energie hörbar summend aus diesem Emitter.
\end{pexample}

    {\small Nahkampf}\\

\begin{tabular}{@{}ll@{}}
    \textbf{\trans{weight_label}:} & 1.00 kg \\
    \textbf{\trans{price_label}:} & 5000 \\
    \textbf{Durchschlag} & 1\\
    \textbf{Verborgenheit} & 2\\
    \textbf{Schadenspotential} & 3\\
\end{tabular}

\wbif{phasesix}{
    \vspace{3mm}
\faIcon{space-shuttle}\hspace{0.5em}}{}



\subsection*{Pugio}

\begin{pexample}
Der Pugio ist ein römischer Dolch von ca. 40cm Länge. Er wurde von den Legionären im antiken Rom als Zweitwaffe getragen. Er besitzt eine breite, blattförmige Stahlklinge, die oft in der Mitte tailliert ist, und einen einfachen, meist T-förmigen Griff aus Metall, Knochen oder Horn.
\end{pexample}

    {\small Nahkampf}\\

\begin{tabular}{@{}ll@{}}
    \textbf{\trans{weight_label}:} & 0.50 kg \\
    \textbf{\trans{price_label}:} & 120 \\
    \textbf{Schadenspotential} & 1\\
\end{tabular}

\wbif{phasesix}{
    \vspace{3mm}
\faIcon{eye}\hspace{0.5em}}{}



\subsection*{Machete}

\begin{pexample}
Sende keine Botschaft!
\end{pexample}

    {\small Nahkampf}\\

\begin{tabular}{@{}ll@{}}
    \textbf{\trans{weight_label}:} & 1.20 kg \\
    \textbf{\trans{price_label}:} & 450 \\
    \textbf{Verborgenheit} & 4\\
    \textbf{Schadenspotential} & 3\\
\end{tabular}

\wbif{phasesix}{
    \vspace{3mm}
\faIcon{chess-rook}\hspace{0.5em}\faIcon{umbrella}\hspace{0.5em}\faIcon{flag}\hspace{0.5em}\faIcon{plane}\hspace{0.5em}\faIcon{microchip}\hspace{0.5em}}{}



\subsection*{Flammendolch}

\begin{pexample}
Die Klinge dieses Dolches ist nicht glatt, sondern in einem wellenförmigen, „flammenartigen“ Muster geschmiedet. Der Flammendolch hat eine gezackte Klinge, die dafür bekannt ist, tiefe, reißende Wunden zu schlagen und Bluten zu verursachen. Der Griff ist meist einfach gehalten, oft aus dunklem Metall oder Leder, um den Fokus auf den unregelmäßigen Stahl zu legen.
\end{pexample}

    {\small Nahkampf}\\

\begin{tabular}{@{}ll@{}}
    \textbf{\trans{weight_label}:} & 0.80 kg \\
    \textbf{\trans{price_label}:} & 450 \\
    \textbf{Verborgenheit} & 2\\
    \textbf{Vorbereitung} & 0\\
    \textbf{Schadenspotential} & 1\\
    \textbf{Blutend} & 1\\
\end{tabular}

\wbif{phasesix}{
    \vspace{3mm}
\faIcon{eye}\hspace{0.5em}\faIcon{chess-rook}\hspace{0.5em}\faIcon{umbrella}\hspace{0.5em}\faIcon{flag}\hspace{0.5em}\faIcon{plane}\hspace{0.5em}\faIcon{microchip}\hspace{0.5em}}{}



\subsection*{Panzerbrecher}

\begin{pexample}
Ein schwerer Dolch, welcher zum Beschädigen von Rüstungen geeignet ist. Für jeden Treffer verliert das Opfer einen Schutz nach seiner Wahl zusätzlich zu evtl. eingesetztem Schutz.
\end{pexample}

    {\small Nahkampf}\\

\begin{tabular}{@{}ll@{}}
    \textbf{\trans{weight_label}:} & 2.00 kg \\
    \textbf{\trans{price_label}:} & 650 \\
    \textbf{Verborgenheit} & 3\\
    \textbf{Schadenspotential} & 1\\
    \textbf{Panzerbrechend} & 1\\
\end{tabular}

\wbif{phasesix}{
    \vspace{3mm}
\faIcon{chess-rook}\hspace{0.5em}\faIcon{umbrella}\hspace{0.5em}\faIcon{flag}\hspace{0.5em}\faIcon{plane}\hspace{0.5em}\faIcon{microchip}\hspace{0.5em}}{}



\subsection*{Degen}

\begin{pexample}
Oft als die Waffe eines Adligen angesehen, wird dieses leichte Schwert durch seine schlanke, steife Klinge definiert, die typischerweise einen dreieckigen Querschnitt aufweist und zu einer außergewöhnlich feinen Spitze zuläuft. Ihr fehlt eine echte Schneide, da sie eine reine Stichwaffe ist. Der Handschutz ist ihr markantestes Merkmal und besteht aus zwei kleinen Schalen und einem zierlichen Knöchelschutz, der auf Präzision ausgelegt ist.
\end{pexample}

    {\small Nahkampf}\\

\begin{tabular}{@{}ll@{}}
    \textbf{\trans{weight_label}:} & 1.20 kg \\
    \textbf{\trans{price_label}:} & 140 \\
    \textbf{Durchschlag} & 1\\
    \textbf{Verborgenheit} & 3\\
    \textbf{Schadenspotential} & 1\\
\end{tabular}

\wbif{phasesix}{
    \vspace{3mm}
\faIcon{chess-rook}\hspace{0.5em}\faIcon{umbrella}\hspace{0.5em}\faIcon{flag}\hspace{0.5em}\faIcon{plane}\hspace{0.5em}\faIcon{microchip}\hspace{0.5em}}{}



\subsection*{Jagdmesser}

\begin{pexample}
Dieses Messer mit feststehender Klinge besitzt eine kräftige, einschneidige Stahlklinge, oft mit einer leichten Krümmung oder einer Clip-Point-Spitze. Der Griff ist praktisch, aus poliertem Holz, Knochen oder Geweih gefertigt und für einen soliden, sicheren Halt geformt. Es ist ein robustes, funktionales Werkzeug, dem die schlanke Form oder die Doppelschneide eines Kampfdolches fehlt.
\end{pexample}

    {\small Nahkampf}\\

\begin{tabular}{@{}ll@{}}
    \textbf{\trans{weight_label}:} & 0.30 kg \\
    \textbf{\trans{price_label}:} & 120 \\
    \textbf{Durchschlag} & 1\\
    \textbf{Vorbereitung} & 0\\
    \textbf{Schadenspotential} & 1\\
\end{tabular}

\wbif{phasesix}{
    \vspace{3mm}
\faIcon{cube}\hspace{0.5em}}{}



\subsection*{Flammberge}

\begin{pexample}
Dieses massive Zweihandschwert ist sofort an seiner Klinge erkennbar. Der lange, schwere, zweischneidige Stahl ist nicht gerade, sondern von der Parierstange bis zur Spitze in einer durchgehenden, wellenförmigen oder „flammenartigen“ Form geschmiedet. Sein Griff ist für eine kraftvolle zweihändige Führung ausgelegt und weist oft eine komplexe Parierstange mit vorwärts gerichteten Haken oder Seitenringen auf.
\end{pexample}

    {\small Nahkampf}\\

\begin{tabular}{@{}ll@{}}
    \textbf{\trans{weight_label}:} & 1.40 kg \\
    \textbf{\trans{price_label}:} & 2200 \\
    \textbf{Reichweite} & 2\\
    \textbf{Verborgenheit} & 6\\
    \textbf{Schadenspotential} & 3\\
    \textbf{Blutend} & 1\\
\end{tabular}

\wbif{phasesix}{
    \vspace{3mm}
\faIcon{eye}\hspace{0.5em}\faIcon{chess-rook}\hspace{0.5em}\faIcon{umbrella}\hspace{0.5em}\faIcon{flag}\hspace{0.5em}\faIcon{plane}\hspace{0.5em}\faIcon{microchip}\hspace{0.5em}}{}



\subsection*{Kettensäge}

\begin{pexample}
Ein klobiges, motorisiertes Gehäuse (oft leuchtend rot oder gelb lackiert) bildet den Hauptkörper, ausgestattet mit einem hinteren Pistolengriff und einem oberen Haltegriff. Von der Vorderseite ragt eine lange, flache Metall-Führungsschiene hervor. Eine endlose Kette, besetzt mit zahlreichen scharfen Metallzähnen, ist um diese Schiene geschlungen und dafür ausgelegt, vom Motormechanismus im Gehäuse mit hoher Geschwindigkeit angetrieben zu werden.
\end{pexample}

    {\small Nahkampf}\\

\begin{tabular}{@{}ll@{}}
    \textbf{\trans{weight_label}:} & 4.00 kg \\
    \textbf{\trans{price_label}:} & 800 \\
    \textbf{Verborgenheit} & 6\\
    \textbf{Vorbereitung} & 2\\
    \textbf{Kapaziät} & 6\\
    \textbf{Schadenspotential} & 3\\
    \textbf{Durchschlag} & 1\\
\end{tabular}

\wbif{phasesix}{
    \vspace{3mm}
\faIcon{plane}\hspace{0.5em}\faIcon{microchip}\hspace{0.5em}}{}



\subsection*{Wakizashi}

\begin{pexample}
Diese Waffe ist eine einschneidige, gekrümmmte Klinge, merklich kürzer als ein volles Langschwert oder Katana, aber deutlich länger als ein Standarddolch. Sie verfügt über einen Griff, der lang genug für ein oder zwei Hände ist, oft mit Rochenhaut und Kordel umwickelt, und ist mit einem runden oder quadratischen Metall-Stichblatt ausgestattet, das Klinge und Griff trennt.
\end{pexample}

    {\small Nahkampf}\\

\begin{tabular}{@{}ll@{}}
    \textbf{\trans{weight_label}:} & 0.90 kg \\
    \textbf{\trans{price_label}:} & 180 \\
    \textbf{Verborgenheit} & 2\\
    \textbf{Vorbereitung} & 0\\
    \textbf{Schadenspotential} & 2\\
\end{tabular}

\wbif{phasesix}{
    \vspace{3mm}
\faIcon{chess-rook}\hspace{0.5em}\faIcon{umbrella}\hspace{0.5em}\faIcon{flag}\hspace{0.5em}\faIcon{plane}\hspace{0.5em}\faIcon{microchip}\hspace{0.5em}}{}



\subsection*{Sichel}

\begin{pexample}
Die Sichel besteht aus einem kurzen, einfachen Holzgriff für eine einzelne Hand. Daran befestigt ist eine Metallklinge mit einer ausgeprägten, C-förmigen Krümmung, die nur entlang ihrer konkaven Innenseite geschärft ist. Die gesamte Konstruktion ist rudimentär und ähnelt einem gewöhnlichen Erntegerät.
\end{pexample}

    {\small Nahkampf}\\

\begin{tabular}{@{}ll@{}}
    \textbf{\trans{weight_label}:} & 0.80 kg \\
    \textbf{\trans{price_label}:} & 150 \\
    \textbf{Verborgenheit} & 4\\
    \textbf{Schadenspotential} & 2\\
\end{tabular}

\wbif{phasesix}{
    \vspace{3mm}
\faIcon{eye}\hspace{0.5em}\faIcon{chess-rook}\hspace{0.5em}\faIcon{umbrella}\hspace{0.5em}\faIcon{flag}\hspace{0.5em}\faIcon{plane}\hspace{0.5em}\faIcon{microchip}\hspace{0.5em}}{}



\subsection*{Rapier}

\begin{pexample}
Eine lange, außergewöhnlich schlanke und steife Klinge, die zu einer nadelspitzen Spitze zuläuft, definiert dieses einhändige Schwert. Sein markantestes Merkmal ist das Gefäß, das einen komplizierten und komplexen Handschutz aus ineinandergreifenden Metallstangen, Ringen oder einer Glocke besitzt, der die Hand des Trägers vollständig schützen soll.
\end{pexample}

    {\small Nahkampf}\\

\begin{tabular}{@{}ll@{}}
    \textbf{\trans{weight_label}:} & 1.00 kg \\
    \textbf{\trans{price_label}:} & 1200 \\
    \textbf{Verborgenheit} & 3\\
    \textbf{Schadenspotential} & 3\\
\end{tabular}

\wbif{phasesix}{
    \vspace{3mm}
\faIcon{cube}\hspace{0.5em}}{}



\subsection*{Breitschwert}

\begin{pexample}
Dieses Schwert zeichnet sich durch eine deutlich breite, gerade, zweischneidige Klinge aus, die sich nur leicht zu einer funktionalen Spitze verjüngt. Oft als brutale Waffe angesehen, sind seine schwere Bauweise und die einfache, robuste Parierstange klar auf rohe Gewalt ausgelegt. Die gesamte Waffe liegt schwer in der Hand und lässt die Raffinesse eines Rapiers oder die Balance eines Ritterschwerts vermissen.
\end{pexample}

    {\small Nahkampf}\\

\begin{tabular}{@{}ll@{}}
    \textbf{\trans{weight_label}:} & 3.00 kg \\
    \textbf{\trans{price_label}:} & 500 \\
    \textbf{Verborgenheit} & 5\\
    \textbf{Schadenspotential} & 4\\
\end{tabular}

\wbif{phasesix}{
    \vspace{3mm}
\faIcon{chess-rook}\hspace{0.5em}\faIcon{umbrella}\hspace{0.5em}\faIcon{flag}\hspace{0.5em}\faIcon{plane}\hspace{0.5em}\faIcon{microchip}\hspace{0.5em}}{}



\subsection*{Molekularschwert}

\begin{pexample}
Ein Griff aus dunklem, nicht-reflektierendem Verbundmaterial bildet das Heft und weist einen einfachen Aktivierungsknopf auf. Daraus ragt ein dünner, starrer Grat hervor, der eine Schneide stützt, die so fein ist, dass sie beinahe unsichtbar wirkt und im direkten Licht leicht schimmert. Die Klinge ist vollkommen gerade und besitzt nicht die Dicke oder den Schliff eines traditionellen Schwertes.
\end{pexample}

    {\small Nahkampf}\\

\begin{tabular}{@{}ll@{}}
    \textbf{\trans{weight_label}:} & 0.80 kg \\
    \textbf{\trans{price_label}:} & 3000 \\
    \textbf{Durchschlag} & 2\\
    \textbf{Verborgenheit} & 5\\
    \textbf{Schadenspotential} & 3\\
\end{tabular}

\wbif{phasesix}{
    \vspace{3mm}
\faIcon{space-shuttle}\hspace{0.5em}}{}



\subsection*{Linkhand}

\begin{pexample}
Die Linkhand ist eine Verteidigungswaffe. In der linken Hand geführt, bietet sie eine zusätzliche Aktion pro Kampfrunde, die nur zur Verteidigung genutzt werden kann.
\end{pexample}

    {\small Nahkampf}\\

\begin{tabular}{@{}ll@{}}
    \textbf{\trans{weight_label}:} & 0.60 kg \\
    \textbf{\trans{price_label}:} & 580 \\
    \textbf{Verborgenheit} & 2\\
    \textbf{Vorbereitung} & 0\\
    \textbf{Schadenspotential} & 1\\
\end{tabular}

\wbif{phasesix}{
    \vspace{3mm}
\faIcon{chess-rook}\hspace{0.5em}\faIcon{umbrella}\hspace{0.5em}\faIcon{flag}\hspace{0.5em}\faIcon{plane}\hspace{0.5em}\faIcon{microchip}\hspace{0.5em}}{}



\subsection*{Zweililien}

\begin{pexample}
Das Zweililien besteht aus einem etwa eineinhalb Schritt langen Stab, an dessen Ende jeweils eine Klinge eingearbeitet ist. Die Waffe wird in der Regel beidhändig geführt.
\end{pexample}

    {\small Nahkampf}\\

\begin{tabular}{@{}ll@{}}
    \textbf{\trans{weight_label}:} & 1.10 kg \\
    \textbf{\trans{price_label}:} & 950 \\
    \textbf{Reichweite} & 2\\
    \textbf{Verborgenheit} & 7\\
    \textbf{Schadenspotential} & 2\\
\end{tabular}

\wbif{phasesix}{
    \vspace{3mm}
\faIcon{chess-rook}\hspace{0.5em}\faIcon{umbrella}\hspace{0.5em}\faIcon{flag}\hspace{0.5em}\faIcon{plane}\hspace{0.5em}\faIcon{microchip}\hspace{0.5em}}{}



\subsection*{Langschwert}

\begin{pexample}
Eine gerade, zweischneidige Klinge, die zu einer scharfen, funktionalen Spitze zuläuft, bildet den Kern dieser Waffe. Sie ist mit einer einfachen, geraden Metall-Parierstange und einem für eine einzelne Hand bemessenen Griff ausgestattet, der von einem schweren Metallknauf ausbalanciert wird. Dies ist das klassische kreuzförmige Schwert, länger als ein Dolch, aber kürzer als eine echte Zweihandwaffe.
\end{pexample}

    {\small Nahkampf}\\

\begin{tabular}{@{}ll@{}}
    \textbf{\trans{weight_label}:} & 1.40 kg \\
    \textbf{\trans{price_label}:} & 2000 \\
    \textbf{Reichweite} & 2\\
    \textbf{Verborgenheit} & 5\\
    \textbf{Schadenspotential} & 2\\
\end{tabular}

\wbif{phasesix}{
    \vspace{3mm}
\faIcon{chess-rook}\hspace{0.5em}}{}



\subsection*{Stilett}

\begin{pexample}
Ein Parierdolch, wenn dieser Dolch geführt wird, erhält der Träger eine zusätzliche Aktion pro Kampfrunde, die nur für die Verteidigung genutzt werden darf.
\end{pexample}

    {\small Nahkampf}\\

\begin{tabular}{@{}ll@{}}
    \textbf{\trans{weight_label}:} & 0.80 kg \\
    \textbf{\trans{price_label}:} & 750 \\
    \textbf{Verborgenheit} & 2\\
    \textbf{Vorbereitung} & 0\\
    \textbf{Schadenspotential} & 1\\
\end{tabular}

\wbif{phasesix}{
    \vspace{3mm}
\faIcon{chess-rook}\hspace{0.5em}\faIcon{umbrella}\hspace{0.5em}\faIcon{flag}\hspace{0.5em}\faIcon{plane}\hspace{0.5em}\faIcon{microchip}\hspace{0.5em}}{}



\subsection*{Dolch}

\begin{pexample}
Ein einfacher Dolch, geschmiedet aus Eisen. Dieser Gegenstand wird sowohl als Werkzeug als auch als Waffe werwendet.
\end{pexample}

    {\small Nahkampf}\\

\begin{tabular}{@{}ll@{}}
    \textbf{\trans{weight_label}:} & 0.60 kg \\
    \textbf{\trans{price_label}:} & 300 \\
    \textbf{Verborgenheit} & 1\\
    \textbf{Vorbereitung} & 0\\
    \textbf{Schadenspotential} & 1\\
\end{tabular}

\wbif{phasesix}{
    \vspace{3mm}
\faIcon{cube}\hspace{0.5em}}{}



\subsection*{Taschenmesser}

\begin{pexample}
Dieser Gegenstand ist ein kompaktes Handwerkzeug mit einem Gehäuse aus Holz, Metall oder Kunststoff. Eine oder mehrere Metallklingen sind in den Griff eingeklappt. Diese Klingen können aus dem Gehäuse herausgeschwenkt werden.
\end{pexample}

    {\small Nahkampf}\\

\begin{tabular}{@{}ll@{}}
    \textbf{\trans{weight_label}:} & 0.10 kg \\
    \textbf{\trans{price_label}:} & 20 \\
    \textbf{Vorbereitung} & 0\\
\end{tabular}

\wbif{phasesix}{
    \vspace{3mm}
\faIcon{flag}\hspace{0.5em}\faIcon{plane}\hspace{0.5em}\faIcon{microchip}\hspace{0.5em}}{}



\subsection*{Pata}

\begin{pexample}
Diese Waffe kombiniert einen starren Stahlhandschuh, der Hand und Unterarm schützt, mit einer langen, geraden, zweischneidigen Klinge. Die Klinge ragt direkt aus dem Ende des Handschuhs hervor, ausgerichtet mit dem Arm des Trägers. Der Griff ist ein horizontaler Querbalken, der im Inneren des geschlossenen Handschuhs gehalten wird und die Hand vollständig verbirgt.
\end{pexample}

    {\small Nahkampf}\\

\begin{tabular}{@{}ll@{}}
    \textbf{\trans{weight_label}:} & 1.00 kg \\
    \textbf{\trans{price_label}:} & 850 \\
    \textbf{Durchschlag} & 1\\
    \textbf{Verborgenheit} & 3\\
    \textbf{Schadenspotential} & 3\\
\end{tabular}

\wbif{phasesix}{
    \vspace{3mm}
\faIcon{chess-rook}\hspace{0.5em}\faIcon{umbrella}\hspace{0.5em}\faIcon{flag}\hspace{0.5em}\faIcon{plane}\hspace{0.5em}\faIcon{microchip}\hspace{0.5em}}{}



\subsection*{Klappspaten}

\begin{pexample}
Ein kurzer, robuster Stiel aus Holz oder Metall ist an einem viereckigen, spitzen Schaufelblatt angebracht. Das hervorstechendste Merkmal des Kopfes ist ein massives Scharnier und ein Feststellring an seiner Basis. Dieser Mechanismus erlaubt es, das Metallblatt in einem 90-Grad-Winkel oder komplett flach gegen den Stiel einzuklappen.
\end{pexample}

    {\small Nahkampf}\\

\begin{tabular}{@{}ll@{}}
    \textbf{\trans{weight_label}:} & 1.20 kg \\
    \textbf{\trans{price_label}:} & 20 \\
    \textbf{Schadenspotential} & 1\\
    \textbf{Verborgenheit} & 3\\
\end{tabular}

\wbif{phasesix}{
    \vspace{3mm}
\faIcon{flag}\hspace{0.5em}\faIcon{plane}\hspace{0.5em}\faIcon{microchip}\hspace{0.5em}}{}



\subsection*{Gladius}

\begin{pexample}
Der Gladius (Mehrzahl: Gladii) ist ein römisches kurzes (ca.80cm) Schwert. Er soll im späteren 3. Jahrhundert v. Chr. aus einem Schwerttyp der iberischen Keltiberer entwickelt worden sein und war in Variationen bis in das 3. Jahrhundert n. Chr. die Standardwaffe der Infanterie der römischen Armee.
\end{pexample}

    {\small Nahkampf}\\

\begin{tabular}{@{}ll@{}}
    \textbf{\trans{weight_label}:} & 1.00 kg \\
    \textbf{\trans{price_label}:} & 200 \\
    \textbf{Schadenspotential} & 3\\
\end{tabular}

\wbif{phasesix}{
    \vspace{3mm}
\faIcon{eye}\hspace{0.5em}}{}



\section{Äxte}



\subsection*{Handaxt}

\begin{pexample}
Ein kurzer, einhändiger Holzstiel ist mit einem einfachen, keilförmigen Stahlkopf versehen. Dieser Kopf weist eine einzelne, ausgestellte Schneide auf, die durch einen flachen, schmucklosen Nacken auf der Gegenseite ausbalanciert wird. Es ist ein kompaktes, zweckmäßiges Werkzeug, dem der Dorn oder Bart einer reinen Streitaxt fehlt.
\end{pexample}

    {\small Nahkampf}\\

\begin{tabular}{@{}ll@{}}
    \textbf{\trans{weight_label}:} & 0.30 kg \\
    \textbf{\trans{price_label}:} & 80 \\
    \textbf{Vorbereitung} & 0\\
    \textbf{Schadenspotential} & 2\\
\end{tabular}

\wbif{phasesix}{
    \vspace{3mm}
\faIcon{cube}\hspace{0.5em}}{}



\subsection*{Feuerwehraxt}

\begin{pexample}
Ein langer, gerader Stiel, oft leuchtend rot oder gelb lackiert, ist mit einem schweren Stahlkopf bestückt. Dieser Kopf ist sofort erkennbar und weist auf der einen Seite eine Standard-Axtschneide und auf der gegenüberliegenden Seite einen scharfen, panzerbrechenden Dorn oder Haken auf. Es ist ein schweres, funktionales Werkzeug, dem die Balance einer reinen Streitaxt fehlt.
\end{pexample}

    {\small Nahkampf}\\

\begin{tabular}{@{}ll@{}}
    \textbf{\trans{weight_label}:} & 2.00 kg \\
    \textbf{\trans{price_label}:} & 300 \\
    \textbf{Durchschlag} & 1\\
    \textbf{Verborgenheit} & 3\\
    \textbf{Schadenspotential} & 3\\
\end{tabular}

\wbif{phasesix}{
    \vspace{3mm}
\faIcon{umbrella}\hspace{0.5em}\faIcon{flag}\hspace{0.5em}\faIcon{plane}\hspace{0.5em}\faIcon{microchip}\hspace{0.5em}}{}



\subsection*{Kriegsaxt}

\begin{pexample}
Ein schwerer, geschmiedeter Stahlkopf, der für den Kampf ausbalanciert ist, krönt diese Waffe. Eine Seite weist eine breite, ausgestellte Schneide (einen „Bart“) auf. Die der Klinge gegenüberliegende Seite ist zu einem scharfen Dorn ausgezogen, was sie deutlich von einer einfachen Allzweckaxt unterscheidet.
\end{pexample}

    {\small Nahkampf}\\

\begin{tabular}{@{}ll@{}}
    \textbf{\trans{weight_label}:} & 5.00 kg \\
    \textbf{\trans{price_label}:} & 2000 \\
    \textbf{Durchschlag} & 1\\
    \textbf{Reichweite} & 2\\
    \textbf{Verborgenheit} & 8\\
    \textbf{Schadenspotential} & 4\\
\end{tabular}

\wbif{phasesix}{
    \vspace{3mm}
\faIcon{chess-rook}\hspace{0.5em}}{}



\subsection*{Axt}

\begin{pexample}
Ein keilförmiger Kopf aus Gusseisen oder Stahl ist auf einem robusten, einhändigen Holzschaft montiert. Der Kopf besitzt eine einzelne, ausgestellte Schneide auf der einen Seite und einen flachen, schweren Nacken auf der anderen. Das Design ist einfach und vielseitig und besitzt weder die Länge einer Langaxt noch die spezielle Balance einer Wurfaxt.
\end{pexample}

    {\small Nahkampf}\\

\begin{tabular}{@{}ll@{}}
    \textbf{\trans{weight_label}:} & 2.00 kg \\
    \textbf{\trans{price_label}:} & 250 \\
    \textbf{Verborgenheit} & 4\\
    \textbf{Schadenspotential} & 3\\
\end{tabular}

\wbif{phasesix}{
    \vspace{3mm}
\faIcon{chess-rook}\hspace{0.5em}\faIcon{umbrella}\hspace{0.5em}\faIcon{flag}\hspace{0.5em}\faIcon{plane}\hspace{0.5em}\faIcon{microchip}\hspace{0.5em}}{}



\subsection*{Holzfälleraxt}

\begin{pexample}
Ein langer, gerader Stiel aus glattem, schmucklosem Hartholz bietet einen zweihändigen Griff. Er trägt einen schweren, keilförmigen Stahlkopf. Dieser Kopf weist eine einzelne breite, extrem scharfe Schneide auf, die durch einen dicken, flachen Nacken auf der Gegenseite ausbalanciert wird.
\end{pexample}

    {\small Nahkampf}\\

\begin{tabular}{@{}ll@{}}
    \textbf{\trans{weight_label}:} & 2.20 kg \\
    \textbf{\trans{price_label}:} & 350 \\
    \textbf{Durchschlag} & 1\\
    \textbf{Verborgenheit} & 5\\
    \textbf{Schadenspotential} & 3\\
\end{tabular}

\wbif{phasesix}{
    \vspace{3mm}
\faIcon{chess-rook}\hspace{0.5em}\faIcon{umbrella}\hspace{0.5em}\faIcon{flag}\hspace{0.5em}\faIcon{plane}\hspace{0.5em}\faIcon{microchip}\hspace{0.5em}}{}



\subsection*{Langaxt}

\begin{pexample}
Die Langaxt zeichnet sich durch einen langen, stabilen Holzschaft aus, der für einen zweihändigen Griff gebaut und oft mit Metallbändern verstärkt ist. Sie ist mit einem einzelnen, schweren Axtkopf bestückt, der typischerweise eine breite, ausladende Schneide besitzt. Anders als einer Hellebarde fehlt diesem Design eine Speerspitze oder ein rückwärtiger Dorn, wodurch es sich rein auf die Kraft der Hauptklinge konzentriert.
\end{pexample}

    {\small Nahkampf}\\

\begin{tabular}{@{}ll@{}}
    \textbf{\trans{weight_label}:} & 2.50 kg \\
    \textbf{\trans{price_label}:} & 500 \\
    \textbf{Reichweite} & 2\\
    \textbf{Verborgenheit} & 5\\
    \textbf{Schadenspotential} & 3\\
\end{tabular}

\wbif{phasesix}{
    \vspace{3mm}
\faIcon{chess-rook}\hspace{0.5em}\faIcon{umbrella}\hspace{0.5em}\faIcon{flag}\hspace{0.5em}\faIcon{plane}\hspace{0.5em}\faIcon{microchip}\hspace{0.5em}}{}



\subsection*{Francisca}

\begin{pexample}
Die Francisca (auch Franzisca) ist eine besondere Form des Wurfbeils, vor allem in Yadosien im ersten und frühen zweiten Jahrhundert verbreitet.
\end{pexample}

    {\small Einzelschuss}\\

\begin{tabular}{@{}ll@{}}
    \textbf{\trans{weight_label}:} & 0.60 kg \\
    \textbf{\trans{price_label}:} & 60 \\
    \textbf{Durchschlag} & 1\\
    \textbf{Vorbereitung} & 0\\
    \textbf{Kapaziät} & 1\\
    \textbf{Schadenspotential} & 2\\
\end{tabular}

\wbif{phasesix}{
    \vspace{3mm}
\faIcon{chess-rook}\hspace{0.5em}\faIcon{umbrella}\hspace{0.5em}\faIcon{flag}\hspace{0.5em}\faIcon{plane}\hspace{0.5em}\faIcon{microchip}\hspace{0.5em}}{}



\subsection*{Tomahawk}

\begin{pexample}
Der Tomahawk  besteht aus einem geraden Schaft, meist aus Holz, der als Griff dient. Am oberen Ende des Schafts ist ein Kopf, typischerweise aus Metall, befestigt, der eine einzelne, vertikal zum Griff ausgerichtete scharfe Klinge aufweist.
\end{pexample}

    {\small Einzelschuss}\\

\begin{tabular}{@{}ll@{}}
    \textbf{\trans{weight_label}:} & 1.20 kg \\
    \textbf{\trans{price_label}:} & 40 \\
    \textbf{Durchschlag} & 1\\
    \textbf{Reichweite} & 15\\
    \textbf{Verborgenheit} & 2\\
    \textbf{Vorbereitung} & 0\\
    \textbf{Kapaziät} & 1\\
    \textbf{Schadenspotential} & 3\\
\end{tabular}

\wbif{phasesix}{
    \vspace{3mm}
\faIcon{chess-rook}\hspace{0.5em}\faIcon{umbrella}\hspace{0.5em}\faIcon{flag}\hspace{0.5em}\faIcon{plane}\hspace{0.5em}\faIcon{microchip}\hspace{0.5em}}{}



\section{Hiebwaffen}



\subsection*{Brechstange}

\begin{pexample}
Eine handelsübliche Brechstange aus massivem, schwerem Stahl, oft sechseckig oder rund im Querschnitt. Ein Ende ist zu einer deutlichen L- oder S-Form gebogen und läuft in einer abgeflachten, gespaltenen Spitze aus. Das gegenüberliegende Ende ist typischerweise zu einem einfachen, keilförmigen Meißel abgeflacht.
\end{pexample}

    {\small Nahkampf}\\

\begin{tabular}{@{}ll@{}}
    \textbf{\trans{weight_label}:} & 3.00 kg \\
    \textbf{\trans{price_label}:} & 40 \\
    \textbf{Schadenspotential} & 1\\
\end{tabular}

\wbif{phasesix}{
    \vspace{3mm}
\faIcon{flag}\hspace{0.5em}\faIcon{plane}\hspace{0.5em}\faIcon{microchip}\hspace{0.5em}\faIcon{space-shuttle}\hspace{0.5em}}{}



\subsection*{Treibhammer}

\begin{pexample}
Ein enormer, zylindrischer Kopf, oft aus dichtem Hartholz gefertigt und mit dicken Eisenringen gebunden, sitzt auf einem dicken, zweihändigen Holzschaft. Anders als die flachen Metallflächen eines Vorschlaghammers oder der Dorn eines Kriegshammers, ist diese Waffe reine, brutale, stumpfe Masse. Das gesamte Objekt ist schwer und kopflastig, eindeutig für zerschmetternde Einschläge konzipiert.
\end{pexample}

    {\small Nahkampf}\\

\begin{tabular}{@{}ll@{}}
    \textbf{\trans{weight_label}:} & 3.00 kg \\
    \textbf{\trans{price_label}:} & 800 \\
    \textbf{Verborgenheit} & 5\\
    \textbf{Schadenspotential} & 4\\
\end{tabular}

\wbif{phasesix}{
    \vspace{3mm}
\faIcon{chess-rook}\hspace{0.5em}}{}



\subsection*{Streitkolben}

\begin{pexample}
Diese Waffe verfügt über einen soliden Schaft aus Holz oder Metall, der für einen einhändigen Griff ausgelegt ist. Der Kopf besteht aus einem schweren, soliden Metallgewicht, das oft mit hervorstehenden Flanschen, Knäufen oder pyramidenförmigen Spitzen gegossen ist. Anders als ein Hammer besitzt diese Waffe keine flache Schlagfläche, sondern konzentriert ihre gesamte Masse auf diese Vorsprünge.
\end{pexample}

    {\small Nahkampf}\\

\begin{tabular}{@{}ll@{}}
    \textbf{\trans{weight_label}:} & 2.80 kg \\
    \textbf{\trans{price_label}:} & 1200 \\
    \textbf{Verborgenheit} & 5\\
    \textbf{Schadenspotential} & 3\\
\end{tabular}

\wbif{phasesix}{
    \vspace{3mm}
\faIcon{chess-rook}\hspace{0.5em}\faIcon{umbrella}\hspace{0.5em}\faIcon{flag}\hspace{0.5em}\faIcon{plane}\hspace{0.5em}\faIcon{microchip}\hspace{0.5em}}{}



\subsection*{Kriegsflegel}

\begin{pexample}
Ein robuster Holzschaft, für eine Hand bemessen, ist über mehrere Glieder einer schweren Kette mit einem frei schwingenden Metallkopf verbunden. Dieser Kopf ist eine massive Eisenkugel, die mit zahlreichen fixierten, scharfen Stacheln besetzt ist. Diese flexible Verbindung unterscheidet ihn unmittelbar vom starren, fixierten Kopf eines Streitkolbens oder Morgensterns.
\end{pexample}

    {\small Nahkampf}\\

\begin{tabular}{@{}ll@{}}
    \textbf{\trans{weight_label}:} & 3.00 kg \\
    \textbf{\trans{price_label}:} & 580 \\
    \textbf{Durchschlag} & 2\\
    \textbf{Reichweite} & 2\\
    \textbf{Verborgenheit} & 6\\
    \textbf{Schadenspotential} & 3\\
\end{tabular}

\wbif{phasesix}{
    \vspace{3mm}
\faIcon{chess-rook}\hspace{0.5em}\faIcon{umbrella}\hspace{0.5em}\faIcon{flag}\hspace{0.5em}\faIcon{plane}\hspace{0.5em}\faIcon{microchip}\hspace{0.5em}}{}



\subsection*{Rabenschnabel}

\begin{pexample}
Diese lange Hiebwaffe ist mit einem metallenen Hammerkopf ausgestattet, der gerne auch in der Form eines Rabenkopfes gefertigt ist. Eine Seite besteht aus einer stumpfen Hammerfläche, während die Gegenseite zu einem langen, scharfen Dorn ausgeformt ist, der dem namensgebenden Rabenschnabel ähnelt. Diese Konstruktion sitzt auf einem langen Holzschaft.
\end{pexample}

    {\small Nahkampf}\\

\begin{tabular}{@{}ll@{}}
    \textbf{\trans{weight_label}:} & 4.00 kg \\
    \textbf{\trans{price_label}:} & 1800 \\
    \textbf{Reichweite} & 2\\
    \textbf{Verborgenheit} & 6\\
    \textbf{Schadenspotential} & 3\\
\end{tabular}

\wbif{phasesix}{
    \vspace{3mm}
\faIcon{chess-rook}\hspace{0.5em}\faIcon{umbrella}\hspace{0.5em}\faIcon{flag}\hspace{0.5em}\faIcon{plane}\hspace{0.5em}\faIcon{microchip}\hspace{0.5em}}{}



\subsection*{Dreschflegel}

\begin{pexample}
Der Dreschflegel besteht aus zwei Holzteilen - einem langen Stab, der als Griff dient, und einem kürzeren, frei schwingenden Schlagstück. Beide Teile sind an einem Ende durch ein lockeres Scharnier aus Leder oder eine kurze Kette verbunden. Es handelt sich sichtlich um ein zweckentfremdetes Erntewerkzeug, dem die Metallstacheln oder ausbalancierten Gewichte eines militärischen Kriegsflegels fehlen.
\end{pexample}

    {\small Nahkampf}\\

\begin{tabular}{@{}ll@{}}
    \textbf{\trans{weight_label}:} & 2.50 kg \\
    \textbf{\trans{price_label}:} & 120 \\
    \textbf{Reichweite} & 2\\
    \textbf{Verborgenheit} & 5\\
    \textbf{Schadenspotential} & 1\\
\end{tabular}

\wbif{phasesix}{
    \vspace{3mm}
\faIcon{eye}\hspace{0.5em}\faIcon{chess-rook}\hspace{0.5em}\faIcon{umbrella}\hspace{0.5em}\faIcon{flag}\hspace{0.5em}\faIcon{plane}\hspace{0.5em}\faIcon{microchip}\hspace{0.5em}}{}



\subsection*{Morgenstern}

\begin{pexample}
Diese Waffe verfügt über einen robusten Schaft aus Holz oder Metall, der für eine Hand ausbalanciert ist und von einem massiven, schweren Kopf gekrönt wird. Dieser oft kugelförmige Kopf ist mit zahlreichen scharfen, feststehenden Stacheln besetzt, was ihn von den stumpfen Klingenflanschen eines Streitkolbens unterscheidet. Das Design ist eindeutig darauf ausgelegt, wuchtige Einschläge mit panzerbrechenden Spitzen zu kombinieren.
\end{pexample}

    {\small Nahkampf}\\

\begin{tabular}{@{}ll@{}}
    \textbf{\trans{weight_label}:} & 3.00 kg \\
    \textbf{\trans{price_label}:} & 1600 \\
    \textbf{Durchschlag} & 1\\
    \textbf{Verborgenheit} & 5\\
    \textbf{Schadenspotential} & 3\\
\end{tabular}

\wbif{phasesix}{
    \vspace{3mm}
\faIcon{chess-rook}\hspace{0.5em}\faIcon{umbrella}\hspace{0.5em}\faIcon{flag}\hspace{0.5em}\faIcon{plane}\hspace{0.5em}\faIcon{microchip}\hspace{0.5em}}{}



\subsection*{Keule}

\begin{pexample}
Die Keule ist ein einzelnes, massives Stück behauenes Hartholz, das durch seine grobe, kopflastige Balance definiert wird. Sie schwillt von einem rauen Griff, der für eine Hand bemessen ist, zu einem dicken, wuchtigen Schlagende an. Ihr fehlen jegliche Metallkomponenten, Stacheln oder Flansche, wodurch sie weit simpler und primitiver als ein Streitkolben wirkt.
\end{pexample}

    {\small Nahkampf}\\

\begin{tabular}{@{}ll@{}}
    \textbf{\trans{weight_label}:} & 1.20 kg \\
    \textbf{\trans{price_label}:} & 15 \\
    \textbf{Verborgenheit} & 4\\
    \textbf{Schadenspotential} & 1\\
\end{tabular}

\wbif{phasesix}{
    \vspace{3mm}
\faIcon{eye}\hspace{0.5em}\faIcon{chess-rook}\hspace{0.5em}\faIcon{umbrella}\hspace{0.5em}\faIcon{flag}\hspace{0.5em}\faIcon{plane}\hspace{0.5em}\faIcon{microchip}\hspace{0.5em}}{}



\subsection*{Schlagstock}

\begin{pexample}
Diese Waffe ist ein solider, gerader Zylinder aus poliertem Hartholz, dichtem Polymer oder schwarzem Stahl, bemessen für einen einhändigen Griff. Ein Ende ist oft für einen sicheren Halt texturiert oder umwickelt, manchmal mit einer ledernen Handschlaufe versehen. Das gesamte Objekt hat eine einheitliche Dicke, ohne den beschwerten Kopf eines Streitkolbens oder einer Keule.
\end{pexample}

    {\small Nahkampf}\\

\begin{tabular}{@{}ll@{}}
    \textbf{\trans{weight_label}:} & 0.80 kg \\
    \textbf{\trans{price_label}:} & 800 \\
    \textbf{Verborgenheit} & 3\\
    \textbf{Schadenspotential} & 1\\
\end{tabular}

\wbif{phasesix}{
    \vspace{3mm}
\faIcon{cube}\hspace{0.5em}}{}



\subsection*{Schmiedehammer}

\begin{pexample}
Ein kurzer, dicker Stiel aus dichtem Hartholz stützt einen schweren, asymmetrischen Block aus Schmiedestahl. Eine Seite dieses Kopfes ist eine breite, flache, quadratische Fläche, die stark von Schlägen und Ruß gezeichnet ist. Die Gegenseite verjüngt sich zu einem stumpfen Keil oder einer abgerundeten „Finne“, wodurch sie sich deutlich von den scharfen Spitzen eines Kriegshammers oder den Zwillingsflächen eines Vorschlaghammers unterscheidet.
\end{pexample}

    {\small Nahkampf}\\

\begin{tabular}{@{}ll@{}}
    \textbf{\trans{weight_label}:} & 4.00 kg \\
    \textbf{\trans{price_label}:} & 850 \\
    \textbf{Verborgenheit} & 4\\
    \textbf{Schadenspotential} & 3\\
\end{tabular}

\wbif{phasesix}{
    \vspace{3mm}
\faIcon{chess-rook}\hspace{0.5em}\faIcon{umbrella}\hspace{0.5em}\faIcon{flag}\hspace{0.5em}\faIcon{plane}\hspace{0.5em}\faIcon{microchip}\hspace{0.5em}}{}



\subsection*{Metallbeschlagene Keule}

\begin{pexample}
Ein dickes, schweres Stück grob behauenes Holz dient als Kern dieser Waffe. Das Schlagende ist dicht mit Eisennieten, stumpfen Bolzen oder geschärften Nagelköpfen besetzt, die tief in die Maserung getrieben wurden. Die Waffe ist kopflastig und grob gefertigt; ihr fehlt die Ausgewogenheit oder die Metallflansche eines militärischen Streitkolbens.
\end{pexample}

    {\small Nahkampf}\\

\begin{tabular}{@{}ll@{}}
    \textbf{\trans{weight_label}:} & 1.50 kg \\
    \textbf{\trans{price_label}:} & 120 \\
    \textbf{Verborgenheit} & 5\\
    \textbf{Schadenspotential} & 2\\
\end{tabular}

\wbif{phasesix}{
    \vspace{3mm}
\faIcon{eye}\hspace{0.5em}\faIcon{chess-rook}\hspace{0.5em}\faIcon{umbrella}\hspace{0.5em}\faIcon{flag}\hspace{0.5em}\faIcon{plane}\hspace{0.5em}\faIcon{microchip}\hspace{0.5em}}{}



\subsection*{Haarkamm}

\begin{pexample}
Ein aus speziell gehärtetem Stahl gefertigter Haarkamm zum Einstecken in Frisuren. Oben eine große Rose, die als Griff dient. Die 5 Zacken des Kamms dienen dann als Klingen. Eine teuflische versteckte Waffe.
\end{pexample}

    {\small Nahkampf}\\

\begin{tabular}{@{}ll@{}}
    \textbf{\trans{weight_label}:} & 0.00 kg \\
    \textbf{\trans{price_label}:} & 100 \\
    \textbf{Verborgenheit} & 1\\
    \textbf{Vorbereitung} & 0\\
    \textbf{Schadenspotential} & 1\\
\end{tabular}

\wbif{phasesix}{
    \vspace{3mm}
\faIcon{microchip}\hspace{0.5em}}{}



\subsection*{Spitzhacke}

\begin{pexample}
Die Spitzhacke ist eigentlich als Handwerkszeug gedacht und besteht aus einem langen, robusten Holzschaft, der für einen zweihändigen Griff ausgelegt ist. Oben ist ein schwerer, geschmiedeter Metallkopf quer montiert. Dieser weist auf der einen Seite einen langen, scharfen Dorn und auf der gegenüberliegenden Seite eine schmale, horizontale Meißelklinge auf. Dadurch bildet er zusammen mit dem Stiel eine T-Form.
\end{pexample}

    {\small Nahkampf}\\

\begin{tabular}{@{}ll@{}}
    \textbf{\trans{weight_label}:} & 2.00 kg \\
    \textbf{\trans{price_label}:} & 180 \\
    \textbf{Durchschlag} & 1\\
    \textbf{Verborgenheit} & 5\\
    \textbf{Schadenspotential} & 2\\
\end{tabular}

\wbif{phasesix}{
    \vspace{3mm}
\faIcon{chess-rook}\hspace{0.5em}\faIcon{umbrella}\hspace{0.5em}\faIcon{flag}\hspace{0.5em}\faIcon{plane}\hspace{0.5em}\faIcon{microchip}\hspace{0.5em}}{}



\subsection*{Vorschlaghammer}

\begin{pexample}
Der Vorschlaghammer besteht aus einem außergewöhnlich langen, dicken Holzschaft, der einen zweihändigen Griff erfordert. Er ist mit einem enormen, schweren Kopf aus einem massiven Block Schmiedeeisens oder Stahls bestückt. Der Kopf weist zwei breite, flache Schlagflächen auf, wobei das Design eindeutig schieres Gewicht und Wucht über die Stacheln oder Flansche eines Kriegshammers stellt.
\end{pexample}

    {\small Nahkampf}\\

\begin{tabular}{@{}ll@{}}
    \textbf{\trans{weight_label}:} & 4.00 kg \\
    \textbf{\trans{price_label}:} & 850 \\
    \textbf{Durchschlag} & 1\\
    \textbf{Reichweite} & 2\\
    \textbf{Verborgenheit} & 6\\
    \textbf{Schadenspotential} & 2\\
\end{tabular}

\wbif{phasesix}{
    \vspace{3mm}
\faIcon{chess-rook}\hspace{0.5em}\faIcon{umbrella}\hspace{0.5em}\faIcon{flag}\hspace{0.5em}\faIcon{plane}\hspace{0.5em}\faIcon{microchip}\hspace{0.5em}}{}



\subsection*{Nunchaku}

\begin{pexample}
Diese Waffe besteht aus zwei kurzen, massiven Hartholzstäben, die an ihren Enden durch eine kurze Metallkette oder ein Seil verbunden sind. Ihre flexible Konstruktion macht sie unberechenbar; erzielt der Trefferwurf keinen Erfolg, so wird der Träger der Waffe für eine Wunde verwundet.
\end{pexample}

    {\small Nahkampf}\\

\begin{tabular}{@{}ll@{}}
    \textbf{\trans{weight_label}:} & 1.20 kg \\
    \textbf{\trans{price_label}:} & 650 \\
    \textbf{Verborgenheit} & 2\\
    \textbf{Vorbereitung} & 0\\
    \textbf{Schadenspotential} & 3\\
\end{tabular}

\wbif{phasesix}{
    \vspace{3mm}
\faIcon{chess-rook}\hspace{0.5em}\faIcon{umbrella}\hspace{0.5em}\faIcon{flag}\hspace{0.5em}\faIcon{plane}\hspace{0.5em}\faIcon{microchip}\hspace{0.5em}}{}



\subsection*{Sturmsense}

\begin{pexample}
Die Kriegssense besteht aus einem langen, verstärkten Holzschaft, ähnlich einer Stangenwaffe. Eine einzelne, große Klinge, klar erkennbar als ein umgeschmiedetes landwirtschaftliches Sensenblatt, ist oben in einer Tülle befestigt und neu ausgerichtet, sodass sie geradeaus in einer Linie mit dem Schaft hervorragt. Die Klinge selbst ist lang und besitzt eine tiefe, konkave Krümmung, wobei ihre geschärfte Schneide auf der Innenseite liegt.
\end{pexample}

    {\small Nahkampf}\\

\begin{tabular}{@{}ll@{}}
    \textbf{\trans{weight_label}:} & 3.00 kg \\
    \textbf{\trans{price_label}:} & 1300 \\
    \textbf{Durchschlag} & 1\\
    \textbf{Reichweite} & 2\\
    \textbf{Verborgenheit} & 7\\
    \textbf{Schadenspotential} & 3\\
\end{tabular}

\wbif{phasesix}{
    \vspace{3mm}
\faIcon{chess-rook}\hspace{0.5em}\faIcon{umbrella}\hspace{0.5em}\faIcon{flag}\hspace{0.5em}\faIcon{plane}\hspace{0.5em}\faIcon{microchip}\hspace{0.5em}}{}



\subsection*{Peitsche}

\begin{pexample}
Die Peitsche ist ein einzelnes, durchgehendes Stück aus dunklem, geflochtenem Leder, oft mehrere Meter lang. Es beginnt mit einem kurzen, starren, geflochtenen Griff, der nahtlos in einen langen, flexiblen Riemen übergeht. Dieser Riemen verjüngt sich gleichmäßig vom dicken Griff bis hinab zu einer sehr feinen, dünnen Kordel an der Spitze.
\end{pexample}

    {\small Nahkampf}\\

\begin{tabular}{@{}ll@{}}
    \textbf{\trans{weight_label}:} & 1.00 kg \\
    \textbf{\trans{price_label}:} & 20 \\
    \textbf{Reichweite} & 3\\
    \textbf{Verborgenheit} & 1\\
    \textbf{Kapaziät} & 1\\
    \textbf{Schadenspotential} & 1\\
\end{tabular}

\wbif{phasesix}{
    \vspace{3mm}
\faIcon{chess-rook}\hspace{0.5em}\faIcon{umbrella}\hspace{0.5em}\faIcon{flag}\hspace{0.5em}\faIcon{plane}\hspace{0.5em}\faIcon{microchip}\hspace{0.5em}}{}



\subsection*{Kriegshammer}

\begin{pexample}
Ein schwerer, geschmiedeter Metallkopf definiert diese Waffe und unterscheidet sie klar von einem Streitkolben. Eine Seite dieses Kopfes weist eine flache oder leicht konvexe, stumpfe Schlagfläche auf, während die Gegenseite zu einem dicken, scharfen Dorn oder einer leicht gebogenen, panzerbrechenden Finne ausgezogen ist. Dieser Kopf ist auf einem robusten, oft für einen Zweihandgriff bemessenen Holz- oder stahlverstärkten Schaft montiert.
\end{pexample}

    {\small Nahkampf}\\

\begin{tabular}{@{}ll@{}}
    \textbf{\trans{weight_label}:} & 5.00 kg \\
    \textbf{\trans{price_label}:} & 1500 \\
    \textbf{Reichweite} & 2\\
    \textbf{Verborgenheit} & 6\\
    \textbf{Vorbereitung} & 2\\
    \textbf{Schadenspotential} & 2\\
\end{tabular}

\wbif{phasesix}{
    \vspace{3mm}
\faIcon{chess-rook}\hspace{0.5em}}{}



\subsection*{Neunschwänzige}

\begin{pexample}
Ein kurzer, solider Griff, oft mit dunklem Leder umwickelt, dient als Handhabe für diese Peitsche. Neun einzelne, gleichlange Riemen oder Kordeln sind an diesem Griff befestigt, von denen jeder in einem harten Knoten oder einem kleinen Metallwiderhaken endet. Die Waffe ist unhandlich: Wenn der Angriffswurf misslingt, fügt die Neunschwänzige dem Träger eine Wunde zu.
\end{pexample}

    {\small Nahkampf}\\

\begin{tabular}{@{}ll@{}}
    \textbf{\trans{weight_label}:} & 1.50 kg \\
    \textbf{\trans{price_label}:} & 850 \\
    \textbf{Reichweite} & 2\\
    \textbf{Verborgenheit} & 5\\
    \textbf{Schadenspotential} & 3\\
\end{tabular}

\wbif{phasesix}{
    \vspace{3mm}
\faIcon{eye}\hspace{0.5em}\faIcon{chess-rook}\hspace{0.5em}\faIcon{umbrella}\hspace{0.5em}\faIcon{flag}\hspace{0.5em}\faIcon{plane}\hspace{0.5em}\faIcon{microchip}\hspace{0.5em}}{}



\subsection*{Schlagring}

\begin{pexample}
Ein einzelnes, solides Stück Gussmetall, typischerweise Messing oder Stahl, ist so geformt, dass es um die Finger passt. Es weist vier runde Löcher auf, durch welche die Finger gesteckt werden. Die Außenkante, die über den Knöcheln sitzt, ist eine dicke, beschwerte Schlagfläche, während ein glatter, gebogener Steg an der Handfläche anliegt.
\end{pexample}

    {\small Nahkampf}\\

\begin{tabular}{@{}ll@{}}
    \textbf{\trans{weight_label}:} & 0.40 kg \\
    \textbf{\trans{price_label}:} & 300 \\
    \textbf{Verborgenheit} & 1\\
    \textbf{Schadenspotential} & 2\\
\end{tabular}

\wbif{phasesix}{
    \vspace{3mm}
\faIcon{flag}\hspace{0.5em}\faIcon{plane}\hspace{0.5em}\faIcon{microchip}\hspace{0.5em}}{}



\subsection*{Eve}


    {\small Nahkampf}\\

\begin{tabular}{@{}ll@{}}
    \textbf{\trans{weight_label}:} & 2.00 kg \\
    \textbf{\trans{price_label}:} & 200 \\
    \textbf{Schadenspotential} & 2\\
    \textbf{Verborgenheit} & 9\\
    \textbf{Reichweite} & 1\\
\end{tabular}

\wbif{phasesix}{
    \vspace{3mm}
\faIcon{umbrella}\hspace{0.5em}\faIcon{industry}\hspace{0.5em}\faIcon{ankh}\hspace{0.5em}\faIcon{syringe}\hspace{0.5em}\faIcon{spider}\hspace{0.5em}}{}



\section{Stangenwaffen}



\subsection*{Lanze}

\begin{pexample}
Ein langer, dicker Schaft aus dichtem Holz bildet den Kern dieser schweren Stangenwaffe. Er wird von einer scharfen, konischen oder blattförmigen Stahlspitze gekrönt, die rein zum Durchbohren konzipiert ist. Ein runder stählerner Handschutz (Brechscheibe) ist oft direkt über dem Griff am Schaft angebracht, was sie von einer einfacheren Pike oder einem Speer unterscheidet.
\end{pexample}

    {\small Nahkampf}\\

\begin{tabular}{@{}ll@{}}
    \textbf{\trans{weight_label}:} & 3.00 kg \\
    \textbf{\trans{price_label}:} & 1000 \\
    \textbf{Durchschlag} & 2\\
    \textbf{Reichweite} & 2\\
    \textbf{Verborgenheit} & 8\\
    \textbf{Schadenspotential} & 2\\
\end{tabular}

\wbif{phasesix}{
    \vspace{3mm}
\faIcon{chess-rook}\hspace{0.5em}}{}



\subsection*{Glefe}

\begin{pexample}
Eine lange, einschneidige Klinge, die an ein großes Messer oder ein Kurzschwert erinnert, ist vertikal am Ende eines langen Holzschafts montiert. Diese Klinge ist bündig mit dem Schaft befestigt, wodurch eine Stangenwaffe entsteht, der die komplexen Speerspitzen oder Seitenklingen einer Partisane oder Corseque fehlen. Die Waffe ist für lange, ausholende Hiebe ausbalanciert.
\end{pexample}

    {\small Nahkampf}\\

\begin{tabular}{@{}ll@{}}
    \textbf{\trans{weight_label}:} & 2.80 kg \\
    \textbf{\trans{price_label}:} & 2750 \\
    \textbf{Durchschlag} & 1\\
    \textbf{Reichweite} & 2\\
    \textbf{Verborgenheit} & 6\\
    \textbf{Schadenspotential} & 3\\
\end{tabular}

\wbif{phasesix}{
    \vspace{3mm}
\faIcon{chess-rook}\hspace{0.5em}\faIcon{umbrella}\hspace{0.5em}\faIcon{flag}\hspace{0.5em}\faIcon{plane}\hspace{0.5em}\faIcon{microchip}\hspace{0.5em}}{}



\subsection*{Holzpflöcke}

\begin{pexample}
Ein Holzpflock, gefertigt aus einem angespitzten Ast oder Stamm. In der Handhabung etwas gewöhnungsbedürftig, aber er wirkt Wunder gegen Vampire, wenn er richtig angewendet wird.

Wenn diese Waffe gegen Vampire mit der Coup de grâce Regel angewendet wird, wird die Anzahl der Würfel um die Kraft des Angreifers erhöht.
\end{pexample}

    {\small Nahkampf}\\

\begin{tabular}{@{}ll@{}}
    \textbf{\trans{weight_label}:} & 0.50 kg \\
    \textbf{\trans{price_label}:} & 10 \\
    \textbf{Schadenspotential} & 1\\
\end{tabular}

\wbif{phasesix}{
    \vspace{3mm}
\faIcon{spider}\hspace{0.5em}}{}



\subsection*{Stab}

\begin{pexample}
Ein langes, solides Stück Hartholz, oft knorrig oder glatt poliert. Es hat etwa die Höhe einer Person und ist möglicherweise mit einem geschnitzten Knauf, einem einfachen Kristall oder schmucklos belassen.
\end{pexample}

    {\small Nahkampf}\\

\begin{tabular}{@{}ll@{}}
    \textbf{\trans{weight_label}:} & 0.80 kg \\
    \textbf{\trans{price_label}:} & 100 \\
    \textbf{Reichweite} & 2\\
    \textbf{Verborgenheit} & 5\\
    \textbf{Vorbereitung} & 0\\
    \textbf{Schadenspotential} & 1\\
\end{tabular}

\wbif{phasesix}{
    \vspace{3mm}
\faIcon{cube}\hspace{0.5em}}{}



\subsection*{Sense}

\begin{pexample}
Ein langer, gebogener Holzschaft (Sensenbaum) ist mit zwei querstehenden Handgriffen ausgestattet. An der Basis dieses Schafts ist eine einzelne, sehr lange, gekrümmte Klinge montiert, die im rechten Winkel absteht, wobei ihre scharfe Kante an der konkaven Seite verläuft. Die gesamte Konstruktion ist unhandlich und eindeutig für weite, schwungvolle Bewegungen ausgelegt.
\end{pexample}

    {\small Nahkampf}\\

\begin{tabular}{@{}ll@{}}
    \textbf{\trans{weight_label}:} & 2.50 kg \\
    \textbf{\trans{price_label}:} & 580 \\
    \textbf{Reichweite} & 2\\
    \textbf{Verborgenheit} & 7\\
    \textbf{Schadenspotential} & 3\\
\end{tabular}

\wbif{phasesix}{
    \vspace{3mm}
\faIcon{eye}\hspace{0.5em}\faIcon{chess-rook}\hspace{0.5em}\faIcon{umbrella}\hspace{0.5em}\faIcon{flag}\hspace{0.5em}\faIcon{plane}\hspace{0.5em}\faIcon{microchip}\hspace{0.5em}}{}



\subsection*{Kampfstab}

\begin{pexample}
Perfekte Balance definiert diese Waffe, einen langen, glatten Schaft aus dichtem Hartholz, der oft an beiden Enden mit gebläutem Metall beschlagen ist. Bekannt als die Waffe eines Mönchs, ist ihre leichte und dennoch stabile Konstruktion für sehr schnelle Angriffe ausgelegt.
\end{pexample}

    {\small Nahkampf}\\

\begin{tabular}{@{}ll@{}}
    \textbf{\trans{weight_label}:} & 0.60 kg \\
    \textbf{\trans{price_label}:} & 150 \\
    \textbf{Reichweite} & 2\\
    \textbf{Verborgenheit} & 6\\
    \textbf{Schadenspotential} & 2\\
\end{tabular}

\wbif{phasesix}{
    \vspace{3mm}
\faIcon{eye}\hspace{0.5em}\faIcon{chess-rook}\hspace{0.5em}\faIcon{umbrella}\hspace{0.5em}\faIcon{flag}\hspace{0.5em}\faIcon{plane}\hspace{0.5em}\faIcon{microchip}\hspace{0.5em}}{}



\subsection*{Hakenspieß}

\begin{pexample}
Der Kopf des Hakenspießes besteht aus einer langen, primären Speerspitze. Von der Basis dieser Spitze ragen zwei kürzere, scharfe Zinken oder Flügel ab, die nach vorn oder leicht nach außen gewinkelt sind. Diese gesamte dreizackige Metallbaugruppe ist auf einem langen Holzschaft montiert.
\end{pexample}

    {\small Nahkampf}\\

\begin{tabular}{@{}ll@{}}
    \textbf{\trans{weight_label}:} & 2.50 kg \\
    \textbf{\trans{price_label}:} & 2450 \\
    \textbf{Reichweite} & 2\\
    \textbf{Verborgenheit} & 5\\
    \textbf{Schadenspotential} & 3\\
\end{tabular}

\wbif{phasesix}{
    \vspace{3mm}
\faIcon{chess-rook}\hspace{0.5em}\faIcon{umbrella}\hspace{0.5em}\faIcon{flag}\hspace{0.5em}\faIcon{plane}\hspace{0.5em}\faIcon{microchip}\hspace{0.5em}}{}



\subsection*{Turnierlanze}

\begin{pexample}
Dies ist ein außergewöhnlich langer, dicker Stangenwaffenschaft, gefertigt aus Holz, das oft mit leuchtenden, heraldischen Farben bemalt ist. Ein konischer Handschutz aus Stahl (Brechscheibe) ist über dem Griff angebracht, und anders als einer Kriegslanze fehlt der Spitze ein geschärfter Punkt. Stattdessen ist sie mit einem stumpfen Metallkopf oder einem Krönel, einem kleinen, kronenartigen Aufsatz, versehen.
\end{pexample}

    {\small Nahkampf}\\

\begin{tabular}{@{}ll@{}}
    \textbf{\trans{weight_label}:} & 2.50 kg \\
    \textbf{\trans{price_label}:} & 1400 \\
    \textbf{Reichweite} & 3\\
    \textbf{Verborgenheit} & 7\\
    \textbf{Vorbereitung} & 2\\
    \textbf{Schadenspotential} & 3\\
\end{tabular}

\wbif{phasesix}{
    \vspace{3mm}
\faIcon{chess-rook}\hspace{0.5em}\faIcon{umbrella}\hspace{0.5em}\faIcon{flag}\hspace{0.5em}\faIcon{plane}\hspace{0.5em}\faIcon{microchip}\hspace{0.5em}}{}



\subsection*{Dreizack}

\begin{pexample}
Ein dreizackiger Metallkopf definiert diese Stangenwaffe. Er besitzt eine lange, zentrale Speerspitze, die auf beiden Seiten von zwei kürzeren, ebenso scharfen Zinken flankiert wird, welche manchmal Widerhaken aufweisen. Diese gesamte Konstruktion ist auf einem langen, robusten Holzschaft montiert.
\end{pexample}

    {\small Nahkampf}\\

\begin{tabular}{@{}ll@{}}
    \textbf{\trans{weight_label}:} & 3.00 kg \\
    \textbf{\trans{price_label}:} & 850 \\
    \textbf{Durchschlag} & 1\\
    \textbf{Reichweite} & 2\\
    \textbf{Verborgenheit} & 5\\
    \textbf{Schadenspotential} & 3\\
\end{tabular}

\wbif{phasesix}{
    \vspace{3mm}
\faIcon{eye}\hspace{0.5em}\faIcon{chess-rook}\hspace{0.5em}\faIcon{umbrella}\hspace{0.5em}\faIcon{flag}\hspace{0.5em}\faIcon{plane}\hspace{0.5em}\faIcon{microchip}\hspace{0.5em}}{}



\subsection*{Pike}

\begin{pexample}
Ein außergewöhnlich langer und starrer Holzschaft, deutlich höher als eine Person, dominiert das Design dieser Stangenwaffe. Er ist mit einer sehr kleinen, scharfen Stahlspitze versehen, die rein funktional ist und nicht die breite Klinge eines Speers besitzt. Die gesamte Waffe ist im Nahbereich unhandlich und ausschließlich darauf ausgelegt, Distanz zu wahren.
\end{pexample}

    {\small Nahkampf}\\

\begin{tabular}{@{}ll@{}}
    \textbf{\trans{weight_label}:} & 3.00 kg \\
    \textbf{\trans{price_label}:} & 1800 \\
    \textbf{Durchschlag} & 2\\
    \textbf{Reichweite} & 2\\
    \textbf{Verborgenheit} & 6\\
    \textbf{Schadenspotential} & 3\\
\end{tabular}

\wbif{phasesix}{
    \vspace{3mm}
\faIcon{chess-rook}\hspace{0.5em}}{}



\subsection*{Mistgabel}

\begin{pexample}
Ein langer, einfacher Holzstiel wird von einem Metallkopf gekrönt. Dieser Kopf teilt sich in zwei, drei oder manchmal vier lange, dünne Metallzinken, die zu einer Spitze geschärft sind und sich leicht nach vorne krümmen. Die gesamte Konstruktion ist die eines einfachen landwirtschaftlichen Werkzeugs, dem die Verstärkung oder Balance eines militärischen Dreizacks fehlt.
\end{pexample}

    {\small Nahkampf}\\

\begin{tabular}{@{}ll@{}}
    \textbf{\trans{weight_label}:} & 2.00 kg \\
    \textbf{\trans{price_label}:} & 5 \\
    \textbf{Durchschlag} & 1\\
    \textbf{Reichweite} & 2\\
    \textbf{Verborgenheit} & 1\\
    \textbf{Nachladeaktionen} & 0\\
    \textbf{Schadenspotential} & 2\\
\end{tabular}

\wbif{phasesix}{
    \vspace{3mm}
\faIcon{cube}\hspace{0.5em}}{}



\subsection*{Speer}

\begin{pexample}
Ein langer, gerader Schaft aus poliertem Hartholz bildet den Hauptkörper dieser Stangenwaffe. Er ist mit einer einzelnen, scharfen Metallspitze versehen, die oft zu einer einfachen Blattform oder einem geschärften Dorn geschmiedet ist. Die Waffe ist auf Reichweite ausbalanciert und besitzt nicht die komplexen Flügel einer Partisane oder das Axtblatt einer Hellebarde.
\end{pexample}

    {\small Nahkampf}\\

\begin{tabular}{@{}ll@{}}
    \textbf{\trans{weight_label}:} & 1.20 kg \\
    \textbf{\trans{price_label}:} & 600 \\
    \textbf{Durchschlag} & 1\\
    \textbf{Reichweite} & 2\\
    \textbf{Verborgenheit} & 6\\
    \textbf{Schadenspotential} & 1\\
\end{tabular}

\wbif{phasesix}{
    \vspace{3mm}
\faIcon{cube}\hspace{0.5em}}{}



\subsection*{Hellebarde}

\begin{pexample}
Ein komplexer, vielseitiger Kopf aus geschmiedetem Stahl krönt diese lange Stangenwaffe. Er kombiniert drei verschiedene Elemente: eine nach vorne gerichtete Speerspitze, ein breites, schweres Axtblatt auf der einen Seite und einen scharfen Dorn oder Haken auf der Rückseite. Diese gesamte Baugruppe ist auf einem langen, robusten Holzschaft montiert, der oft mit Metall-Schienen (Langetten) verstärkt ist.
\end{pexample}

    {\small Nahkampf}\\

\begin{tabular}{@{}ll@{}}
    \textbf{\trans{weight_label}:} & 3.00 kg \\
    \textbf{\trans{price_label}:} & 3000 \\
    \textbf{Durchschlag} & 2\\
    \textbf{Reichweite} & 2\\
    \textbf{Verborgenheit} & 8\\
    \textbf{Schadenspotential} & 3\\
\end{tabular}

\wbif{phasesix}{
    \vspace{3mm}
\faIcon{chess-rook}\hspace{0.5em}}{}



\subsection*{Kriegslanze}

\begin{pexample}
Ein langer, dicker Schaft aus dichtem, verstärktem Holz bildet den Kern dieser schweren Stangenwaffe. Er wird von einer scharfen, blattförmigen oder konischen Stahlspitze gekrönt, die zum Durchbohren konzipiert ist und nicht das stumpfe Krönel seines Turnier-Gegenstücks besitzt. Ein runder stählerner Handschutz (Brechscheibe) ist oft direkt über dem Griffbereich am Schaft angebracht.
\end{pexample}

    {\small Nahkampf}\\

\begin{tabular}{@{}ll@{}}
    \textbf{\trans{weight_label}:} & 3.50 kg \\
    \textbf{\trans{price_label}:} & 1500 \\
    \textbf{Durchschlag} & 2\\
    \textbf{Reichweite} & 3\\
    \textbf{Verborgenheit} & 7\\
    \textbf{Schadenspotential} & 3\\
\end{tabular}

\wbif{phasesix}{
    \vspace{3mm}
\faIcon{chess-rook}\hspace{0.5em}\faIcon{umbrella}\hspace{0.5em}\faIcon{flag}\hspace{0.5em}\faIcon{plane}\hspace{0.5em}\faIcon{microchip}\hspace{0.5em}}{}



\subsection*{Partisane}

\begin{pexample}
Ein langer, robuster Holzschaft wird von einer breiten, zweischneidigen Speerspitze gekrönt. An der Basis dieser Hauptklinge flankieren zwei kleinere, symmetrische und geschärfte Vorsprünge, die oft wie Mondsicheln oder spitze Flügel geformt sind. Der gesamte Kopf ist typischerweise flach und kunstvoller als ein Militärspeer, ihm fehlt jedoch das Axtblatt einer Hellebarde.
\end{pexample}

    {\small Nahkampf}\\

\begin{tabular}{@{}ll@{}}
    \textbf{\trans{weight_label}:} & 3.00 kg \\
    \textbf{\trans{price_label}:} & 1000 \\
    \textbf{Durchschlag} & 1\\
    \textbf{Reichweite} & 2\\
    \textbf{Verborgenheit} & 6\\
    \textbf{Schadenspotential} & 3\\
\end{tabular}

\wbif{phasesix}{
    \vspace{3mm}
\faIcon{chess-rook}\hspace{0.5em}\faIcon{umbrella}\hspace{0.5em}\faIcon{flag}\hspace{0.5em}\faIcon{plane}\hspace{0.5em}\faIcon{microchip}\hspace{0.5em}}{}



\section{Bögen}



\subsection*{Kompositbogen}

\begin{pexample}
Dieser Bogen ist aus mehreren laminierten Materialien gefertigt – typischerweise Horn, Sehne und ein Holzkern – die fest mit Wicklungen verbunden sind. Im ungespannten Zustand krümmen sich seine Wurfarme an den Spitzen deutlich nach vorne. Die gesamte Konstruktion ist oft kompakter als ein Schlachtbogen, weist aber eine weitaus komplexere und wohlüberlegte Form auf.
\end{pexample}

    {\small Einzelschuss}\\

\begin{tabular}{@{}ll@{}}
    \textbf{\trans{weight_label}:} & 1.20 kg \\
    \textbf{\trans{price_label}:} & 750 \\
    \textbf{Reichweite} & 50\\
    \textbf{Verborgenheit} & 5\\
    \textbf{Vorbereitung} & 0\\
    \textbf{Kapaziät} & 1\\
    \textbf{Schadenspotential} & 3\\
\end{tabular}

\wbif{phasesix}{
    \vspace{3mm}
\faIcon{chess-rook}\hspace{0.5em}\faIcon{umbrella}\hspace{0.5em}\faIcon{flag}\hspace{0.5em}\faIcon{plane}\hspace{0.5em}\faIcon{microchip}\hspace{0.5em}}{}



\subsection*{Repetierarmbrust}

\begin{pexample}
Ein markanter Hebelmechanismus, der oft in den Schaft integriert und mit einem oben liegenden Kasten- oder Schwerkraftmagazin verbunden ist, definiert diese Armbrust. Es ist eine sehr einfach zu spannende Armbrust, welche schnell zu laden ist. Die restliche Konstruktion besteht aus einem Standard-Wurfarm und Abzug, auch wenn der Schaft stark modifiziert ist, um den Ladevorgang aufzunehmen.
\end{pexample}

    {\small Einzelschuss}\\

\begin{tabular}{@{}ll@{}}
    \textbf{\trans{weight_label}:} & 3.00 kg \\
    \textbf{\trans{price_label}:} & 950 \\
    \textbf{Durchschlag} & 1\\
    \textbf{Reichweite} & 30\\
    \textbf{Verborgenheit} & 3\\
    \textbf{Vorbereitung} & 2\\
    \textbf{Kapaziät} & 2\\
    \textbf{Schadenspotential} & 3\\
\end{tabular}

\wbif{phasesix}{
    \vspace{3mm}
\faIcon{chess-rook}\hspace{0.5em}\faIcon{umbrella}\hspace{0.5em}\faIcon{flag}\hspace{0.5em}\faIcon{plane}\hspace{0.5em}\faIcon{microchip}\hspace{0.5em}}{}



\subsection*{Armbrust}

\begin{pexample}
Ein solider Schaft aus Holz oder Metall bildet den Hauptkörper dieser Waffe und weist auf seiner Oberseite eine Führungsrille auf. Nahe der Front ist ein einzelner, horizontaler Wurfarm montiert. Ein Abzugsmechanismus ist in den Schaft eingelassen, dafür ausgelegt, eine gespannte Bogensehne unter Spannung zu halten.
\end{pexample}

    {\small Einzelschuss}\\

\begin{tabular}{@{}ll@{}}
    \textbf{\trans{weight_label}:} & 3.00 kg \\
    \textbf{\trans{price_label}:} & 800 \\
    \textbf{Durchschlag} & 1\\
    \textbf{Reichweite} & 40\\
    \textbf{Verborgenheit} & 3\\
    \textbf{Vorbereitung} & 2\\
    \textbf{Kapaziät} & 1\\
    \textbf{Nachladeaktionen} & 2\\
    \textbf{Schadenspotential} & 3\\
\end{tabular}

\wbif{phasesix}{
    \vspace{3mm}
\faIcon{eye}\hspace{0.5em}\faIcon{chess-rook}\hspace{0.5em}\faIcon{umbrella}\hspace{0.5em}\faIcon{flag}\hspace{0.5em}}{}



\subsection*{Jagdbogen}

\begin{pexample}
Diese Waffe besteht aus einem langen, einzelnen Stab aus abgelagertem Holz, oft Eibe oder Esche. Sie hat einen einfachen Ledergriff und eine Sehne aus gedrehter Faser, besitzt jedoch nicht die schwere Verstärkung eines Schlachtbogens.
\end{pexample}

    {\small Einzelschuss}\\

\begin{tabular}{@{}ll@{}}
    \textbf{\trans{weight_label}:} & 1.20 kg \\
    \textbf{\trans{price_label}:} & 650 \\
    \textbf{Durchschlag} & 1\\
    \textbf{Reichweite} & 60\\
    \textbf{Kapaziät} & 1\\
    \textbf{Schadenspotential} & 2\\
\end{tabular}

\wbif{phasesix}{
    \vspace{3mm}
\faIcon{eye}\hspace{0.5em}\faIcon{chess-rook}\hspace{0.5em}\faIcon{umbrella}\hspace{0.5em}\faIcon{flag}\hspace{0.5em}\faIcon{plane}\hspace{0.5em}\faIcon{microchip}\hspace{0.5em}}{}



\subsection*{Langbogen}

\begin{pexample}
Dieser Bogen wird durch seine beeindruckende Höhe definiert, die oft der seines Trägers entspricht oder sie übertrifft. Er besteht aus einem einzigen, langen Stab aus abgelagertem Holz, typischerweise Eibe oder Esche, der glattpoliert ist. Anders als der vielschichtige Kompositbogen stammt seine Kraft rein aus der Länge seiner einfachen, leicht gebogenen Wurfarme und der Spannung seiner dicken Sehne.
\end{pexample}

    {\small Einzelschuss}\\

\begin{tabular}{@{}ll@{}}
    \textbf{\trans{weight_label}:} & 2.00 kg \\
    \textbf{\trans{price_label}:} & 600 \\
    \textbf{Durchschlag} & 1\\
    \textbf{Reichweite} & 60\\
    \textbf{Verborgenheit} & 5\\
    \textbf{Kapaziät} & 1\\
    \textbf{Schadenspotential} & 2\\
\end{tabular}

\wbif{phasesix}{
    \vspace{3mm}
\faIcon{cube}\hspace{0.5em}}{}



\subsection*{Kurzbogen}

\begin{pexample}
Ein einzelner Stab aus abgelagertem Holz, oft Esche oder Ulme, formt diesen kompakten Bogen. Seine Wurfarme sind merklich kurz und schlank, und ihm fehlt die beeindruckende Höhe eines Langbogens oder die schwere Verstärkung eines Schlachtbogens. Die Waffe ist leicht und mit einer einfachen, gedrehten Fasersehne ausgestattet.
\end{pexample}

    {\small Einzelschuss}\\

\begin{tabular}{@{}ll@{}}
    \textbf{\trans{weight_label}:} & 1.20 kg \\
    \textbf{\trans{price_label}:} & 400 \\
    \textbf{Reichweite} & 40\\
    \textbf{Verborgenheit} & 3\\
    \textbf{Kapaziät} & 1\\
    \textbf{Schadenspotential} & 2\\
\end{tabular}

\wbif{phasesix}{
    \vspace{3mm}
\faIcon{cube}\hspace{0.5em}}{}



\subsection*{Zweifacharmbrust}

\begin{pexample}
Ein zentraler Schaft aus Holz oder Metall ist mit zwei horizontalen Wurfarmen bestückt, die entweder nebeneinander oder vertikal gestapelt angebracht sind. Die Waffe verfügt über einen Abzugsmechanismus, der mit zwei parallelen Rillen oder Schienen verbunden ist. Diese Führungen sind dafür ausgelegt, zwei separate Bolzen zu halten, welche durch entsprechende Sehnen unter Spannung gesichert werden.
\end{pexample}

    {\small Einzelschuss}\\

\begin{tabular}{@{}ll@{}}
    \textbf{\trans{weight_label}:} & 3.80 kg \\
    \textbf{\trans{price_label}:} & 2900 \\
    \textbf{Durchschlag} & 1\\
    \textbf{Reichweite} & 45\\
    \textbf{Verborgenheit} & 4\\
    \textbf{Vorbereitung} & 2\\
    \textbf{Kapaziät} & 2\\
    \textbf{Schadenspotential} & 3\\
\end{tabular}

\wbif{phasesix}{
    \vspace{3mm}
\faIcon{chess-rook}\hspace{0.5em}\faIcon{umbrella}\hspace{0.5em}\faIcon{flag}\hspace{0.5em}\faIcon{plane}\hspace{0.5em}\faIcon{microchip}\hspace{0.5em}}{}



\subsection*{Handarmbrust}

\begin{pexample}
Diese Waffe ist eine kleine, verborgene Armbrust, die oft mit einem pistolenartigen Griff gebaut ist. Sie besitzt einen kurzen, horizontalen Wurfarm und einen simplen Mechanismus, der sie einfach zu spannen macht. Die gesamte leichte Konstruktion ist darauf ausgelegt, schnell gezogen werden zu können.
\end{pexample}

    {\small Einzelschuss}\\

\begin{tabular}{@{}ll@{}}
    \textbf{\trans{weight_label}:} & 0.80 kg \\
    \textbf{\trans{price_label}:} & 900 \\
    \textbf{Durchschlag} & 1\\
    \textbf{Reichweite} & 10\\
    \textbf{Verborgenheit} & 1\\
    \textbf{Vorbereitung} & 0\\
    \textbf{Kapaziät} & 1\\
    \textbf{Schadenspotential} & 2\\
\end{tabular}

\wbif{phasesix}{
    \vspace{3mm}
\faIcon{chess-rook}\hspace{0.5em}\faIcon{umbrella}\hspace{0.5em}\faIcon{flag}\hspace{0.5em}\faIcon{plane}\hspace{0.5em}\faIcon{microchip}\hspace{0.5em}}{}



\subsection*{Vierfacharmbrust}

\begin{pexample}
Diese schwere Waffe besitzt einen breiten, verstärkten Schaft, der vier horizontale Wurfarme aufnimmt, die oft gestapelt oder nebeneinander angeordnet sind. Die Oberseite des Schafts weist vier parallele Rillen auf, die jeweils einen Bolzen halten. Ein komplexer Abzugsmechanismus ist mit einem komplizierten Sehnensystem verbunden, das alle vier Bolzen gleichzeitig unter Spannung hält.
\end{pexample}

    {\small Einzelschuss}\\

\begin{tabular}{@{}ll@{}}
    \textbf{\trans{weight_label}:} & 6.00 kg \\
    \textbf{\trans{price_label}:} & 4100 \\
    \textbf{Durchschlag} & 1\\
    \textbf{Reichweite} & 45\\
    \textbf{Rückstosskontrolle} & 1\\
    \textbf{Verborgenheit} & 8\\
    \textbf{Vorbereitung} & 2\\
    \textbf{Kapaziät} & 4\\
    \textbf{Nachladeaktionen} & 4\\
    \textbf{Schadenspotential} & 3\\
\end{tabular}

\wbif{phasesix}{
    \vspace{3mm}
\faIcon{chess-rook}\hspace{0.5em}\faIcon{umbrella}\hspace{0.5em}\faIcon{flag}\hspace{0.5em}\faIcon{plane}\hspace{0.5em}\faIcon{microchip}\hspace{0.5em}}{}



\subsection*{Kriegsbogen}

\begin{pexample}
Dieser schwere Bogen besitzt einen dicken, verstärkten Rahmen, oft aus Schichtholz oder Kompositmaterial, was ihn größer und robuster als einen einfachen Jagdbogen macht. Seine Wurfarme sind steif und stark gekrümmt und mit einer dicken, schweren Bogensehne verbunden. Die gesamte Waffe ist auf Kraft ausgelegt und verzichtet zugunsten eines rein funktionalen Designs auf Verzierungen.
\end{pexample}

    {\small Einzelschuss}\\

\begin{tabular}{@{}ll@{}}
    \textbf{\trans{weight_label}:} & 1.40 kg \\
    \textbf{\trans{price_label}:} & 1450 \\
    \textbf{Durchschlag} & 1\\
    \textbf{Reichweite} & 70\\
    \textbf{Verborgenheit} & 6\\
    \textbf{Kapaziät} & 1\\
    \textbf{Schadenspotential} & 4\\
\end{tabular}

\wbif{phasesix}{
    \vspace{3mm}
\faIcon{chess-rook}\hspace{0.5em}\faIcon{umbrella}\hspace{0.5em}\faIcon{flag}\hspace{0.5em}\faIcon{plane}\hspace{0.5em}\faIcon{microchip}\hspace{0.5em}}{}



\subsection*{Leichte Armbrust}

\begin{pexample}
Ein Schaft aus Holz oder Kompositmaterial bildet den Hauptkörper dieser Waffe. An dessen Vorderseite ist ein einzelner, relativ kurzer horizontaler Wurfarm (Prod) befestigt. Sie verfügt über eine einzelne Bogensehne, einen einfachen Abzugsmechanismus und eine Führungsrille auf der Oberseite des Schafts, die einen einzelnen Bolzen aufnimmt.
\end{pexample}

    {\small Einzelschuss}\\

\begin{tabular}{@{}ll@{}}
    \textbf{\trans{weight_label}:} & 2.20 kg \\
    \textbf{\trans{price_label}:} & 650 \\
    \textbf{Durchschlag} & 1\\
    \textbf{Reichweite} & 40\\
    \textbf{Verborgenheit} & 3\\
    \textbf{Kapaziät} & 1\\
    \textbf{Schadenspotential} & 2\\
\end{tabular}

\wbif{phasesix}{
    \vspace{3mm}
\faIcon{chess-rook}\hspace{0.5em}\faIcon{umbrella}\hspace{0.5em}\faIcon{flag}\hspace{0.5em}\faIcon{plane}\hspace{0.5em}\faIcon{microchip}\hspace{0.5em}}{}



\subsection*{Bogen der Elemente}

\begin{pexample}
Während Sie auf diesen Bogen eingestimmt sind und ihn halten, können Sie eine Minute lang meditieren und sich auf eine der fünf Elementarschadensarten konzentrieren. Am Ende der Minute pulsiert der Bogen hell in der Farbe, die dieser Schadensart entspricht, und seine Ausrichtung ändert sich zu diesem Element.

In diesen Bogen sind aufwändige wirbelnde Muster eingraviert, etwa wirbelnde Wolken, peitschende Wellen und Feuerzungen. Die Ritzen der Motive erstrahlen in einem langsam pulsierenden Licht. Die Farbe dieses Lichts hängt davon ab, auf welche Elementarschadensart der Bogen gerade ausgerichtet ist: Rot für Feuer, Weiß für Kälte, Grün für Säure, Blau für Blitz, Lila für Donner.
\end{pexample}

    {\small Einzelschuss}\\

\begin{tabular}{@{}ll@{}}
    \textbf{\trans{weight_label}:} & 1.00 kg \\
    \textbf{\trans{price_label}:} & 100000 \\
    \textbf{Reichweite} & 50\\
    \textbf{Verborgenheit} & 5\\
    \textbf{Kapaziät} & 10\\
    \textbf{Schadenspotential} & 3\\
\end{tabular}

\wbif{phasesix}{
    \vspace{3mm}
\faIcon{magic}\hspace{0.5em}}{}



\section{Pistolen}



\subsection*{Laserpistole}

\begin{pexample}
Ein kompaktes, kantiges Gehäuse aus weißem oder schwarzem Polymer bildet den Korpus dieser Handfeuerwaffe. Anstelle eines Verschlusses oder eines herkömmlichen Laufs besitzt sie einen kurzen, verstärkten Emitter, der von Fokussierlinsen oder metallischen Rippen umgeben ist. Eine schwere, abnehmbare Energiezelle, die oft schwach glüht, wird in die Basis des Griffs eingeschoben.
\end{pexample}

    {\small Feuerstoß}\\
    {\small Einzelschuss}\\

\begin{tabular}{@{}ll@{}}
    \textbf{\trans{weight_label}:} & 0.80 kg \\
    \textbf{\trans{price_label}:} & 2200 \\
    \textbf{Durchschlag} & 2\\
    \textbf{Reichweite} & 60\\
    \textbf{Rückstosskontrolle} & 2\\
    \textbf{Verborgenheit} & 2\\
    \textbf{Kapaziät} & 20\\
    \textbf{Schadenspotential} & 2\\
\end{tabular}

\wbif{phasesix}{
    \vspace{3mm}
\faIcon{space-shuttle}\hspace{0.5em}}{}



\subsection*{Colt Peacemaker}

\begin{pexample}
Ein polierter Metallrahmen definiert diesen Single-Action-Revolver. Sein markantestes Merkmal ist der rotierende, geriffelte Zylinder, der für einzelne Patronen ausgelegt ist und über eine Ladeklappe auf der rechten Rahmenseite geladen wird. Ein charakteristischer, geschwungener Griff (oft aus Holz oder Elfenbein), ein außenliegender Hahn und ein starres Frontkorn vervollständigen sein ikonisches Profil.
\end{pexample}

    {\small Einzelschuss}\\

\begin{tabular}{@{}ll@{}}
    \textbf{\trans{weight_label}:} & 1.00 kg \\
    \textbf{\trans{price_label}:} & 800 \\
    \textbf{Durchschlag} & 1\\
    \textbf{Reichweite} & 20\\
    \textbf{Verborgenheit} & 3\\
    \textbf{Kapaziät} & 6\\
    \textbf{Schadenspotential} & 2\\
\end{tabular}

\wbif{phasesix}{
    \vspace{3mm}
\faIcon{umbrella}\hspace{0.5em}\faIcon{flag}\hspace{0.5em}\faIcon{plane}\hspace{0.5em}\faIcon{microchip}\hspace{0.5em}}{}



\subsection*{Desert Eagle}

\begin{pexample}
Ein außergewöhnlich großer und schwerer Rahmen verleiht dieser halbautomatischen Handfeuerwaffe ihre unverkennbare, kantige Silhouette. Sie besitzt einen dicken, schwerkalibrigen Lauf (oft mit einem markanten dreieckigen Profil) und einen massiven Verschluss. Der überdimensionierte Griff nimmt ein großes, einreihiges Stangenmagazin auf.
\end{pexample}

    {\small Einzelschuss}\\

\begin{tabular}{@{}ll@{}}
    \textbf{\trans{weight_label}:} & 1.00 kg \\
    \textbf{\trans{price_label}:} & 1200 \\
    \textbf{Durchschlag} & 1\\
    \textbf{Reichweite} & 20\\
    \textbf{Verborgenheit} & 3\\
    \textbf{Kapaziät} & 7\\
    \textbf{Schadenspotential} & 3\\
\end{tabular}

\wbif{phasesix}{
    \vspace{3mm}
\faIcon{plane}\hspace{0.5em}\faIcon{microchip}\hspace{0.5em}}{}



\subsection*{Colt Dragoon}

\begin{pexample}
Ein massiver, schwerer Rahmen und ein langer Lauf (oft achteckig) verleihen diesem Single-Action-Revolver seine imposante Silhouette. Er verfügt über eine rotierende Trommel, die für loses Pulver und Kugel ausgelegt ist, nicht für metallische Patronen. Ein markanter, angelenkter Ladehebel sitzt direkt unter dem Lauf, und ein großer, außenliegender Hahn ruht auf dem Rahmen.
\end{pexample}

    {\small Einzelschuss}\\

\begin{tabular}{@{}ll@{}}
    \textbf{\trans{weight_label}:} & 0.50 kg \\
    \textbf{\trans{price_label}:} & 400 \\
    \textbf{Reichweite} & 20\\
    \textbf{Rückstosskontrolle} & 1\\
    \textbf{Verborgenheit} & 2\\
    \textbf{Kapaziät} & 6\\
    \textbf{Schadenspotential} & 2\\
    \textbf{Nachladeaktionen} & 2\\
\end{tabular}

\wbif{phasesix}{
    \vspace{3mm}
\faIcon{umbrella}\hspace{0.5em}\faIcon{flag}\hspace{0.5em}\faIcon{plane}\hspace{0.5em}\faIcon{microchip}\hspace{0.5em}}{}



\subsection*{HK USP}

\begin{pexample}
Diese moderne Handfeuerwaffe ist um einen kastenförmigen, schwarzen Polymerrahmen und einen schweren, rechteckigen Stahl-Verschluss herum gebaut. Ihr Design ist kantig und funktional, mit einem markanten, eckigen Abzugsbügel und (oft) einem außenliegenden Hahn am hinteren Ende des Verschlusses. Einfache Drei-Punkt-Visiere sitzen obenauf.
\end{pexample}

    {\small Einzelschuss}\\

\begin{tabular}{@{}ll@{}}
    \textbf{\trans{weight_label}:} & 1.00 kg \\
    \textbf{\trans{price_label}:} & 500 \\
    \textbf{Durchschlag} & 1\\
    \textbf{Reichweite} & 20\\
    \textbf{Rückstosskontrolle} & 1\\
    \textbf{Kapaziät} & 8\\
    \textbf{Schadenspotential} & 2\\
\end{tabular}

\wbif{phasesix}{
    \vspace{3mm}
\faIcon{microchip}\hspace{0.5em}}{}



\subsection*{Ruger}

\begin{pexample}
Dieses halbautomatische Gewehr ist sofort an seinem traditionellen, einteiligen Holzschaft erkennbar, der einen einfachen Griff formt und an ein älteres M1-Gewehr erinnert. Es verfügt über einen relativ dünnen Lauf mit einem markanten, geschützten Flügelkorn. Ein gebogenes, rechteckiges Kastenmagazin wird direkt vor dem Abzugsbügel eingerastet.
\end{pexample}

    {\small Einzelschuss}\\

\begin{tabular}{@{}ll@{}}
    \textbf{\trans{weight_label}:} & 0.80 kg \\
    \textbf{\trans{price_label}:} & 800 \\
    \textbf{Durchschlag} & 1\\
    \textbf{Reichweite} & 150\\
    \textbf{Verborgenheit} & 2\\
    \textbf{Kapaziät} & 9\\
    \textbf{Schadenspotential} & 2\\
\end{tabular}

\wbif{phasesix}{
    \vspace{3mm}
\faIcon{plane}\hspace{0.5em}\faIcon{microchip}\hspace{0.5em}}{}



\subsection*{Espingole}

\begin{pexample}
Eine auffallend weite, glockenförmige Mündung aus ausgestelltem Messing oder Eisen definiert diese Feuerwaffe. Dieser kurze, schwere Lauf ist mit einem dicken, karabinerlangen Holzschaft verbunden. Ein markanter Steinschloss-Mechanismus ist seitlich am Gehäuse angebracht und vervollständigt das simple, kopflastige Profil.
\end{pexample}

    {\small Einzelschuss}\\

\begin{tabular}{@{}ll@{}}
    \textbf{\trans{weight_label}:} & 1.20 kg \\
    \textbf{\trans{price_label}:} & 800 \\
    \textbf{Reichweite} & 20\\
    \textbf{Rückstosskontrolle} & 1\\
    \textbf{Verborgenheit} & 3\\
    \textbf{Kapaziät} & 1\\
    \textbf{Schadenspotential} & 1\\
    \textbf{Nachladeaktionen} & 2\\
\end{tabular}

\wbif{phasesix}{
    \vspace{3mm}
\faIcon{umbrella}\hspace{0.5em}}{}



\subsection*{Steinschloßpistole}

\begin{pexample}
Ein schwerer, einzelner Lauf ist mit einem ausgeprägten, geschwungenen Hartholzgriff verbunden. Sein markantestes Merkmal ist der komplexe Metallmechanismus an der Seite: ein C-förmiger Hammer, der ein Stück Feuerstein hält und so ausgerichtet ist, dass er auf eine Stahlplatte schlägt, die sich über einer kleinen Zündpfanne befindet.
\end{pexample}

    {\small Einzelschuss}\\

\begin{tabular}{@{}ll@{}}
    \textbf{\trans{weight_label}:} & 0.80 kg \\
    \textbf{\trans{price_label}:} & 500 \\
    \textbf{Reichweite} & 20\\
    \textbf{Verborgenheit} & 2\\
    \textbf{Kapaziät} & 1\\
    \textbf{Schadenspotential} & 4\\
    \textbf{Nachladeaktionen} & 2\\
\end{tabular}

\wbif{phasesix}{
    \vspace{3mm}
\faIcon{umbrella}\hspace{0.5em}\faIcon{flag}\hspace{0.5em}}{}



\subsection*{Glock (9mm)}

\begin{pexample}
Diese moderne Handfeuerwaffe ist sofort an ihrem kantigen, mattschwarzen Polymerrahmen und dem blockigen, rechteckigen Stahlverschluss erkennbar. Ihr fehlt ein außenliegender Hahn, was dem hinteren Ende des Verschlusses ein glattes, hakenfreies Profil verleiht. Ein markanter, eckiger Abzugsbügel und ein simpler, im Abzug integrierter Sicherungsmechanismus vervollständigen ihr hochfunktionales, schmuckloses Design.
\end{pexample}

    {\small Feuerstoß}\\
    {\small Einzelschuss}\\

\begin{tabular}{@{}ll@{}}
    \textbf{\trans{weight_label}:} & 0.30 kg \\
    \textbf{\trans{price_label}:} & 300 \\
    \textbf{Reichweite} & 80\\
    \textbf{Rückstosskontrolle} & 1\\
    \textbf{Verborgenheit} & 1\\
    \textbf{Kapaziät} & 17\\
    \textbf{Schadenspotential} & 2\\
\end{tabular}

\wbif{phasesix}{
    \vspace{3mm}
\faIcon{plane}\hspace{0.5em}\faIcon{microchip}\hspace{0.5em}}{}



\subsection*{Colt 1911}

\begin{pexample}
Diese halbautomatische Handfeuerwaffe ist durch ihren schlanken Ganzmetallrahmen und einen Verschluss mit hinteren Rillen gekennzeichnet. Sie verfügt über einen markanten Griffwinkel, der von texturierten Schalen bedeckt ist, sowie einen hervorstehenden außenliegenden Hahn. Ein dünnes, einreihiges Stangenmagazin wird in die Basis des Griffs eingeschoben.
\end{pexample}

    {\small Einzelschuss}\\
    {\small Feuerstoß}\\

\begin{tabular}{@{}ll@{}}
    \textbf{\trans{weight_label}:} & 0.70 kg \\
    \textbf{\trans{price_label}:} & 800 \\
    \textbf{Reichweite} & 20\\
    \textbf{Rückstosskontrolle} & 1\\
    \textbf{Verborgenheit} & 2\\
    \textbf{Kapaziät} & 7\\
    \textbf{Schadenspotential} & 2\\
\end{tabular}

\wbif{phasesix}{
    \vspace{3mm}
\faIcon{flag}\hspace{0.5em}\faIcon{plane}\hspace{0.5em}\faIcon{microchip}\hspace{0.5em}}{}



\subsection*{Flechette Pistole}

\begin{pexample}
Ein schlanker, kantiger Rahmen aus dunklen Verbundwerkstoffen definiert diese Handfeuerwaffe. Anstelle eines konventionellen runden Laufs besitzt sie einen kurzen, rechteckigen Abschusskanal, der oft von Stabilisierungsflügeln flankiert wird. Die Munition ist ein langes, dünnes Magazin, häufig aus durchsichtigem Polymer, das horizontal über dem Griff eingeschoben wird und die dicht gepackten, nadelartigen Flechettes im Inneren zeigt.
\end{pexample}

    {\small Feuerstoß}\\

\begin{tabular}{@{}ll@{}}
    \textbf{\trans{weight_label}:} & 5.00 kg \\
    \textbf{\trans{price_label}:} & 1500 \\
    \textbf{Durchschlag} & 1\\
    \textbf{Reichweite} & 15\\
    \textbf{Rückstosskontrolle} & 1\\
    \textbf{Verborgenheit} & 2\\
    \textbf{Kapaziät} & 9\\
    \textbf{Schadenspotential} & 4\\
\end{tabular}

\wbif{phasesix}{
    \vspace{3mm}
\faIcon{space-shuttle}\hspace{0.5em}\faIcon{fire}\hspace{0.5em}}{}



\subsection*{Derringer}

\begin{pexample}
Diese Handfeuerwaffe wird durch ihre extreme Kompaktheit definiert und aufgrund ihrer geringen Größe oft als „Damenwaffe“ betrachtet. Sie verfügt typischerweise über eine sehr kurze Über-Unter-Laufanordnung, die auf einem winzigen Rahmen montiert ist. Der Griff ist minimal.
\end{pexample}

    {\small Einzelschuss}\\
    {\small Feuerstoß}\\

\begin{tabular}{@{}ll@{}}
    \textbf{\trans{weight_label}:} & 0.20 kg \\
    \textbf{\trans{price_label}:} & 200 \\
    \textbf{Reichweite} & 10\\
    \textbf{Rückstosskontrolle} & 1\\
    \textbf{Vorbereitung} & 0\\
    \textbf{Kapaziät} & 2\\
    \textbf{Schadenspotential} & 1\\
\end{tabular}

\wbif{phasesix}{
    \vspace{3mm}
\faIcon{umbrella}\hspace{0.5em}\faIcon{flag}\hspace{0.5em}}{}



\subsection*{Walther P99}

\begin{pexample}
Die Walther P99 (9x19mm Parabellum) ist eine moderne, zuverlässige Selbstladepistole im Kaliber 9x19mm.

Sie wurde aufgrund ihrer ausgeprägten Ergonomie und hohen Sicherheitstechnik bei zahlreichen deutschen Polizeibehörden eingeführt. 

Ausgestattet mit einem robusten Polymergriffstück, das sich dank austauschbarer Griffrücken individuell an verschiedene Handgrößen anpassen lässt, bietet sie dem Nutzer eine intuitive Handhabung. 

Ein technisches Markenzeichen der Waffe ist der sogenannte Anti-Stress-Abzug, der speziell dafür konzipiert wurde, auch unter hoher psychischer Belastung eine kontrollierte und präzise Schussabgabe zu gewährleisten.

Abgerundet wird das durchdachte Design durch ein 15-Schuss-Magazin sowie einen deutlich erkennbaren Ladestandsanzeiger, der jederzeit darüber informiert, ob sich eine Patrone im Patronenlager befindet.
\end{pexample}

    {\small Einzelschuss}\\

\begin{tabular}{@{}ll@{}}
    \textbf{\trans{weight_label}:} & 0.70 kg \\
    \textbf{\trans{price_label}:} & 549 \\
    \textbf{Rückstosskontrolle} & 1\\
    \textbf{Durchschlag} & 1\\
    \textbf{Reichweite} & 20\\
    \textbf{Kapaziät} & 15\\
    \textbf{Schadenspotential} & 2\\
    \textbf{Verborgenheit} & 3\\
\end{tabular}

\wbif{phasesix}{
    \vspace{3mm}
\faIcon{dice-six}\hspace{0.5em}\faIcon{plane}\hspace{0.5em}\faIcon{microchip}\hspace{0.5em}}{}



\section{Sturmgewehre}



\subsection*{AK 47}

\begin{pexample}
Ein stark gekrümmtes, bananenförmiges Kastenmagazin aus Metall definiert dieses Militärgewehr. Es basiert auf einem Gehäuse aus Stahlblech und ist mit einem unverwechselbaren, oft polierten, hölzernen Schaft, Pistolengriff und Handschutz ausgestattet. Ein markantes, getunneltes Korn sitzt nahe der Mündung des dunklen Laufs.
\end{pexample}

    {\small Feuerstoß}\\
    {\small Einzelschuss}\\

\begin{tabular}{@{}ll@{}}
    \textbf{\trans{weight_label}:} & 3.80 kg \\
    \textbf{\trans{price_label}:} & 2500 \\
    \textbf{Durchschlag} & 1\\
    \textbf{Reichweite} & 40\\
    \textbf{Verborgenheit} & 5\\
    \textbf{Kapaziät} & 30\\
    \textbf{Schadenspotential} & 3\\
\end{tabular}

\wbif{phasesix}{
    \vspace{3mm}
\faIcon{plane}\hspace{0.5em}\faIcon{microchip}\hspace{0.5em}}{}



\subsection*{G3}

\begin{pexample}
Das G3 Gewehr ist um ein dunkles Gehäuse aus Stahlblech herum gebaut und mit einem markanten, breiten Handschutz, Pistolengriff und Festschaft aus Polymer ausgestattet. Zu seinen erkennbarsten Merkmalen gehören eine hervorstehende Drehtrommel-Kimme (Diopter) und ein gerades, seitenflaches Kastenmagazin aus Metall.
\end{pexample}

    {\small Feuerstoß}\\
    {\small Einzelschuss}\\

\begin{tabular}{@{}ll@{}}
    \textbf{\trans{weight_label}:} & 4.40 kg \\
    \textbf{\trans{price_label}:} & 3800 \\
    \textbf{Reichweite} & 80\\
    \textbf{Verborgenheit} & 6\\
    \textbf{Kapaziät} & 20\\
    \textbf{Schadenspotential} & 3\\
\end{tabular}

\wbif{phasesix}{
    \vspace{3mm}
\faIcon{plane}\hspace{0.5em}\faIcon{microchip}\hspace{0.5em}}{}



\subsection*{M-16}

\begin{pexample}
Dieses moderne Militärgewehr ist aus einem dunklen, mattierten Metallgehäuse und markanten schwarzen Polymer-Anbauteilen (Schaft, Pistolengriff und Handschutz) gebaut. Seine erkennbarsten Merkmale sind der lange, schlanke Lauf, der in einem hervorstehenden, dreieckigen Kornfuß endet, sowie der feste Tragegriff oben auf dem Gehäuse, der gleichzeitig die Kimme integriert.
\end{pexample}

    {\small Feuerstoß}\\
    {\small Einzelschuss}\\

\begin{tabular}{@{}ll@{}}
    \textbf{\trans{weight_label}:} & 2.90 kg \\
    \textbf{\trans{price_label}:} & 3600 \\
    \textbf{Durchschlag} & 1\\
    \textbf{Reichweite} & 100\\
    \textbf{Verborgenheit} & 6\\
    \textbf{Kapaziät} & 30\\
    \textbf{Schadenspotential} & 3\\
\end{tabular}

\wbif{phasesix}{
    \vspace{3mm}
\faIcon{plane}\hspace{0.5em}\faIcon{microchip}\hspace{0.5em}}{}



\subsection*{Steyr AUG}

\begin{pexample}
Dieses moderne Gewehr wird durch sein kompaktes „Bullpup“-Layout definiert, bei dem der Magazinschacht direkt im Schaft hinter dem Pistolengriff platziert ist. Sein Gehäuse ist eine glatte Polymerhülle, oft olivgrün oder schwarz, die von einem markanten Tragegriff dominiert wird, der direkt in ein integriertes, flaches optisches Visier übergeht. Das Design ist schlank und in sich geschlossen und verzichtet auf die separaten, angeschraubten Teile vieler anderer Militärfeuerwaffen.
\end{pexample}

    {\small Feuerstoß}\\
    {\small Einzelschuss}\\

\begin{tabular}{@{}ll@{}}
    \textbf{\trans{weight_label}:} & 3.60 kg \\
    \textbf{\trans{price_label}:} & 2800 \\
    \textbf{Reichweite} & 300\\
    \textbf{Verborgenheit} & 5\\
    \textbf{Kapaziät} & 30\\
    \textbf{Schadenspotential} & 4\\
\end{tabular}

\wbif{phasesix}{
    \vspace{3mm}
\faIcon{plane}\hspace{0.5em}\faIcon{microchip}\hspace{0.5em}}{}



\subsection*{Sako M95}

\begin{pexample}
Ein markanter, seitlich klappbarer Röhren-Metallschaft definiert dieses moderne Militärgewehr. Sein dunkles Gehäuse aus Stahlblech und das hervorstehende Gasrohr über dem Lauf zeigen seine Abstammung von der AK-Familie. Es ist mit schwarzen oder grünen Polymer-Anbauteilen (Handschutz und Pistolengriff) ausgestattet und wird aus einem gebogenen Polymer-Kastenmagazin gespeist.
\end{pexample}

    {\small Einzelschuss}\\
    {\small Feuerstoß}\\

\begin{tabular}{@{}ll@{}}
    \textbf{\trans{weight_label}:} & 3.50 kg \\
    \textbf{\trans{price_label}:} & 2500 \\
    \textbf{Durchschlag} & 1\\
    \textbf{Reichweite} & 400\\
    \textbf{Verborgenheit} & 5\\
    \textbf{Kapaziät} & 30\\
    \textbf{Schadenspotential} & 3\\
\end{tabular}

\wbif{phasesix}{
    \vspace{3mm}
\faIcon{microchip}\hspace{0.5em}}{}



\subsection*{Flechette Gewehr}

\begin{pexample}
Flechette-Gewehre schießen ein paar hundert Stahlpfeile pro Schuss. Sehr tödlich, mit einer großen Streuung. Aber keine sehr schnelle Waffe. Der Salve kann man ausweichen.
\end{pexample}

    {\small Flechette Salve}\\

\begin{tabular}{@{}ll@{}}
    \textbf{\trans{weight_label}:} & 15.00 kg \\
    \textbf{\trans{price_label}:} & 2000 \\
    \textbf{Durchschlag} & 1\\
    \textbf{Reichweite} & 50\\
    \textbf{Kapaziät} & 18\\
    \textbf{Schadenspotential} & 4\\
\end{tabular}

\wbif{phasesix}{
    \vspace{3mm}
\faIcon{space-shuttle}\hspace{0.5em}}{}



\section{Schleudern}



\subsection*{Schleuder}

\begin{pexample}
Eine lange, doppelt geführte Kordel aus geflochtener Faser oder Sehne weist in ihrer Mitte eine breitere, schalenförmige Ledertasche auf. Die beiden von dieser Tasche ausgehenden Kordeln sind gleich lang, wobei ein Ende oft zu einer einfachen Fingerschlaufe geknotet ist.
\end{pexample}

    {\small Einzelschuss}\\

\begin{tabular}{@{}ll@{}}
    \textbf{\trans{weight_label}:} & 0.10 kg \\
    \textbf{\trans{price_label}:} & 20 \\
    \textbf{Reichweite} & 15\\
    \textbf{Vorbereitung} & 0\\
    \textbf{Kapaziät} & 1\\
    \textbf{Schadenspotential} & 1\\
\end{tabular}

\wbif{phasesix}{
    \vspace{3mm}
\faIcon{cube}\hspace{0.5em}}{}



\subsection*{Pilum mit Amentum}

\begin{pexample}
Ein langer, dünner Eisenschaft, der in einer scharfen, pyramidenförmigen Spitze gipfelt, definiert diesen Wurfspeer. Dieser Metallschaft ist tief in einen kürzeren, beschwerten Holzschaft eingelassen. Nahe dem Schwerpunkt der Waffe ist ein langer Lederriemen (Amentum) fest um den Schaft gewickelt, der zu einer Fingerschlaufe geformt ist.
\end{pexample}

    {\small Einzelschuss}\\

\begin{tabular}{@{}ll@{}}
    \textbf{\trans{weight_label}:} & 2.00 kg \\
    \textbf{\trans{price_label}:} & 1 \\
    \textbf{Durchschlag} & 1\\
    \textbf{Reichweite} & 40\\
    \textbf{Kapaziät} & 1\\
    \textbf{Schadenspotential} & 1\\
\end{tabular}

\wbif{phasesix}{
    \vspace{3mm}
\faIcon{eye}\hspace{0.5em}\faIcon{chess-rook}\hspace{0.5em}}{}



\subsection*{Zwille}

\begin{pexample}
Ein Y-förmiger Rahmen, oft aus einem einzelnen Stück Hartholz geschnitzt oder aus gebogenem Metall geformt, dient als Griff. An den beiden oberen Zinken dieses Rahmens sind zwei starke Gummibänder befestigt. Die gegenüberliegenden Enden dieser Bänder sind durch eine kleine, flexible Tasche, typischerweise aus Leder, verbunden.
\end{pexample}

    {\small Einzelschuss}\\

\begin{tabular}{@{}ll@{}}
    \textbf{\trans{weight_label}:} & 0.50 kg \\
    \textbf{\trans{price_label}:} & 30 \\
    \textbf{Reichweite} & 15\\
    \textbf{Verborgenheit} & 1\\
    \textbf{Vorbereitung} & 0\\
    \textbf{Kapaziät} & 1\\
    \textbf{Schadenspotential} & 2\\
\end{tabular}

\wbif{phasesix}{
    \vspace{3mm}
\faIcon{cube}\hspace{0.5em}}{}



\subsection*{Blasrohr}

\begin{pexample}
Diese Waffe ist eine einfache, lange Röhre aus ausgehöhltem Schilf oder poliertem Holz. Ein Ende ist sichtlich geschnitzt und geglättet, um als Mundstück zu dienen. Das gesamte Objekt ist leicht und vollkommen gerade und besitzt keinerlei Mechanismen oder Sehnen, sondern nur den leeren Kanal, der sich durch seine gesamte Länge zieht.
\end{pexample}

    {\small Einzelschuss}\\

\begin{tabular}{@{}ll@{}}
    \textbf{\trans{weight_label}:} & 0.80 kg \\
    \textbf{\trans{price_label}:} & 120 \\
    \textbf{Durchschlag} & 1\\
    \textbf{Verborgenheit} & 5\\
    \textbf{Vorbereitung} & 0\\
    \textbf{Kapaziät} & 1\\
    \textbf{Reichweite} & 15\\
\end{tabular}

\wbif{phasesix}{
    \vspace{3mm}
\faIcon{eye}\hspace{0.5em}\faIcon{chess-rook}\hspace{0.5em}\faIcon{umbrella}\hspace{0.5em}\faIcon{flag}\hspace{0.5em}\faIcon{plane}\hspace{0.5em}\faIcon{microchip}\hspace{0.5em}}{}



\section{Gewehre}



\subsection*{Muskete}

\begin{pexample}
Ein außergewöhnlich langer, schwerer Stahllauf ist in einen durchgehenden Schaft aus dunklem, poliertem Hartholz eingelassen. Sein markantestes Merkmal ist der große, C-förmige Steinschloss-Mechanismus, der seitlich montiert ist. Die Waffe ist sehr lang, mit einem einfachen, runden Glattrohrlauf, und ein dünner metallener Ladestock ist in einem Kanal direkt darunter eingeschoben.
\end{pexample}

    {\small Einzelschuss}\\

\begin{tabular}{@{}ll@{}}
    \textbf{\trans{weight_label}:} & 2.00 kg \\
    \textbf{\trans{price_label}:} & 800 \\
    \textbf{Durchschlag} & 2\\
    \textbf{Reichweite} & 40\\
    \textbf{Verborgenheit} & 6\\
    \textbf{Kapaziät} & 1\\
    \textbf{Schadenspotential} & 1\\
    \textbf{Nachladeaktionen} & 2\\
\end{tabular}

\wbif{phasesix}{
    \vspace{3mm}
\faIcon{umbrella}\hspace{0.5em}\faIcon{flag}\hspace{0.5em}}{}



\subsection*{Langgewehr}

\begin{pexample}
Ein bemerkenswert langer und schlanker achteckiger Lauf, oft aus dunklem, brüniertem Stahl, gibt dieser Feuerwaffe ihren Namen. Dieser Lauf ist in einen durchgehenden, polierten Hartholzschaft (oft Ahorn) eingelassen. Ihre markantesten Merkmale sind der hervorstehende Steinschloss-Mechanismus an der Seite und die geschwungene, halbmondförmige Schaftkappe aus Metall.
\end{pexample}

    {\small Einzelschuss}\\

\begin{tabular}{@{}ll@{}}
    \textbf{\trans{weight_label}:} & 2.00 kg \\
    \textbf{\trans{price_label}:} & 1500 \\
    \textbf{Durchschlag} & 2\\
    \textbf{Reichweite} & 40\\
    \textbf{Verborgenheit} & 5\\
    \textbf{Kapaziät} & 1\\
    \textbf{Schadenspotential} & 2\\
\end{tabular}

\wbif{phasesix}{
    \vspace{3mm}
\faIcon{flag}\hspace{0.5em}\faIcon{plane}\hspace{0.5em}\faIcon{microchip}\hspace{0.5em}}{}



\subsection*{Springfield Gewehr}

\begin{pexample}
Ein durchgehender, dunkler Hartholzschaft umschließt einen langen, brünierten Stahllauf. Sein markantestes Merkmal ist der massive Metall-Verschlussblock, der oben in das Gehäuse eingelassen ist und zum Laden einer einzelnen, großkalibrigen Patrone nach oben und vorne aufklappt. Ein markanter, außenliegender Hahn ruht direkt hinter diesem Mechanismus.
\end{pexample}

    {\small Einzelschuss}\\

\begin{tabular}{@{}ll@{}}
    \textbf{\trans{weight_label}:} & 2.00 kg \\
    \textbf{\trans{price_label}:} & 1400 \\
    \textbf{Durchschlag} & 2\\
    \textbf{Reichweite} & 80\\
    \textbf{Verborgenheit} & 6\\
    \textbf{Schadenspotential} & 2\\
\end{tabular}

\wbif{phasesix}{
    \vspace{3mm}
\faIcon{flag}\hspace{0.5em}\faIcon{plane}\hspace{0.5em}}{}



\subsection*{Barrett}

\begin{pexample}
Eine enorme, pfeilförmige Mündungsbremse sitzt am Ende eines außergewöhnlich langen, schwerkalibrigen Laufs. Diese massive Feuerwaffe ist um ein schweres, kantiges Stahlgehäuse und einen einfachen Pistolengriff herum konstruiert, was Standard-Militärgewehre dagegen klein wirken lässt. Ein klappbares Zweibein ist nahe der Front montiert, und ein großes, rechteckiges Kastenmagazin wird vor der Abzugseinheit eingeschoben.
\end{pexample}

    {\small Einzelschuss}\\

\begin{tabular}{@{}ll@{}}
    \textbf{\trans{weight_label}:} & 3.00 kg \\
    \textbf{\trans{price_label}:} & 2500 \\
    \textbf{Durchschlag} & 1\\
    \textbf{Reichweite} & 1400\\
    \textbf{Verborgenheit} & 5\\
    \textbf{Kapaziät} & 11\\
    \textbf{Schadenspotential} & 3\\
\end{tabular}

\wbif{phasesix}{
    \vspace{3mm}
\faIcon{plane}\hspace{0.5em}\faIcon{microchip}\hspace{0.5em}}{}



\subsection*{M1}

\begin{pexample}
Ein dunkles, schweres Stahlgehäuse, das in einen einteiligen Vollholzschaft eingelassen ist, verleiht diesem Militärgewehr sein charakteristisches Profil. Es verfügt über ein markantes Lochvisier (Diopter) am hinteren Ende des Gehäuses und ein geschütztes Flügelkorn. Am auffälligsten an seinem Design ist das Fehlen eines abnehmbaren Kastenmagazins; stattdessen nutzt es ein internes Magazin, das von oben mittels eines blockartigen Metall-Laderahmens geladen wird.
\end{pexample}

    {\small Feuerstoß}\\
    {\small Einzelschuss}\\

\begin{tabular}{@{}ll@{}}
    \textbf{\trans{weight_label}:} & 3.00 kg \\
    \textbf{\trans{price_label}:} & 3000 \\
    \textbf{Durchschlag} & 1\\
    \textbf{Reichweite} & 270\\
    \textbf{Verborgenheit} & 5\\
    \textbf{Kapaziät} & 8\\
    \textbf{Schadenspotential} & 4\\
\end{tabular}

\wbif{phasesix}{
    \vspace{3mm}
\faIcon{flag}\hspace{0.5em}\faIcon{plane}\hspace{0.5em}\faIcon{microchip}\hspace{0.5em}}{}



\subsection*{Kar98}

\begin{pexample}
Das Kar basiert auf einem dunklen, einteiligen Holzschaft, was sie zu einem merklich kürzeren und handlicheren Karabiner als das Gewehr 98 macht. Ihre entscheidenden Merkmale sind der metallene Zylinderverschluss-Mechanismus mit einem deutlich nach unten gebogenen Kammerstengel sowie das auf dem Lauf montierte Schiebevisier mit V-Kimme.
\end{pexample}

    {\small Einzelschuss}\\

\begin{tabular}{@{}ll@{}}
    \textbf{\trans{weight_label}:} & 4.00 kg \\
    \textbf{\trans{price_label}:} & 3500 \\
    \textbf{Durchschlag} & 2\\
    \textbf{Reichweite} & 1200\\
    \textbf{Verborgenheit} & 7\\
    \textbf{Kapaziät} & 5\\
    \textbf{Schadenspotential} & 2\\
\end{tabular}

\wbif{phasesix}{
    \vspace{3mm}
\faIcon{flag}\hspace{0.5em}\faIcon{plane}\hspace{0.5em}\faIcon{microchip}\hspace{0.5em}}{}



\subsection*{Winchester '76}

\begin{pexample}
Ein langer, schwerer Lauf, oft achteckig, ist mit einem polierten Hartholzschaft und einem soliden, dunklen Metallgehäuse verbunden. Das erkennbarste Merkmal ist der große Hebelbügel hinter dem Abzugsbügel, der den Mechanismus bedient. Ein außenliegender Hahn ruht auf dem Gehäuse, und eine deutliche Ladeklappe aus Messing oder Stahl ist seitlich in das Gehäuse eingelassen.
\end{pexample}

    {\small Einzelschuss}\\

\begin{tabular}{@{}ll@{}}
    \textbf{\trans{weight_label}:} & 4.00 kg \\
    \textbf{\trans{price_label}:} & 4500 \\
    \textbf{Durchschlag} & 1\\
    \textbf{Reichweite} & 40\\
    \textbf{Verborgenheit} & 5\\
    \textbf{Kapaziät} & 7\\
    \textbf{Schadenspotential} & 3\\
\end{tabular}

\wbif{phasesix}{
    \vspace{3mm}
\faIcon{flag}\hspace{0.5em}\faIcon{plane}\hspace{0.5em}\faIcon{microchip}\hspace{0.5em}}{}



\subsection*{Lasergewehr}

\begin{pexample}
Eine polierte, schwere Kristalllinse an der Mündung zeichnet dieses Gewehr aus. Sein schlanker, kantiger Rahmen besteht aus dunklen Legierungen und Polymeren, dem ein konventioneller Hülsenauswurf oder Magazinschacht fehlt. Stattdessen ist eine klobige, glühende Energiezelle im Schaft arretiert, die Energie durch sichtbare externe Leitungen zuführt.
\end{pexample}

    {\small Feuerstoß}\\
    {\small Einzelschuss}\\

\begin{tabular}{@{}ll@{}}
    \textbf{\trans{weight_label}:} & 1.00 kg \\
    \textbf{\trans{price_label}:} & 5200 \\
    \textbf{Durchschlag} & 2\\
    \textbf{Reichweite} & 800\\
    \textbf{Rückstosskontrolle} & 2\\
    \textbf{Verborgenheit} & 4\\
    \textbf{Kapaziät} & 30\\
    \textbf{Schadenspotential} & 4\\
\end{tabular}

\wbif{phasesix}{
    \vspace{3mm}
\faIcon{space-shuttle}\hspace{0.5em}}{}



\subsection*{Gewehr 98}

\begin{pexample}
Ein langer, durchgehender Holzschaft umschließt einen dunklen Stahllauf, was diese Feuerwaffe deutlich länger als einen Karabiner macht. Ihr entscheidendes Merkmal ist ein metallener Zylinderverschluss-Mechanismus mit einem geraden, polierten Kammerstengel. Ein markantes, auf weite Distanzen ausgelegtes Visier ist auf dem Lauf montiert, und der Schaft ist mit Stahlbändern sowie einer Bajonett-Halterung nahe der Mündung versehen.
\end{pexample}

    {\small Einzelschuss}\\

\begin{tabular}{@{}ll@{}}
    \textbf{\trans{weight_label}:} & 4.00 kg \\
    \textbf{\trans{price_label}:} & 1200 \\
    \textbf{Durchschlag} & 2\\
    \textbf{Reichweite} & 200\\
    \textbf{Verborgenheit} & 5\\
    \textbf{Kapaziät} & 5\\
    \textbf{Schadenspotential} & 2\\
\end{tabular}

\wbif{phasesix}{
    \vspace{3mm}
\faIcon{flag}\hspace{0.5em}\faIcon{plane}\hspace{0.5em}}{}



\section{Maschinenpistolen}



\subsection*{MP5}

\begin{pexample}
Die MP5 basiert auf einem dunklen Gehäuse aus gestanztem Metall und wirkt wie eine verkleinerte Version eines G3-Gewehrs. Sie verfügt über ein markantes, getunneltes Korn und eine Drehtrommel-Kimme (Diopter). Ein gebogenes Kastenmagazin wird direkt vor dem Pistolengriff und der Abzugsgruppe aus Polymer eingeführt, und sie ist oft mit einer einfachen, einschiebbaren Metallschulterstütze ausgestattet.
\end{pexample}

    {\small Feuerstoß}\\
    {\small Einzelschuss}\\

\begin{tabular}{@{}ll@{}}
    \textbf{\trans{weight_label}:} & 3.10 kg \\
    \textbf{\trans{price_label}:} & 1200 \\
    \textbf{Reichweite} & 40\\
    \textbf{Verborgenheit} & 4\\
    \textbf{Kapaziät} & 30\\
    \textbf{Schadenspotential} & 2\\
\end{tabular}

\wbif{phasesix}{
    \vspace{3mm}
\faIcon{plane}\hspace{0.5em}\faIcon{microchip}\hspace{0.5em}}{}



\subsection*{MP40}

\begin{pexample}
Diese automatische Feuerwaffe ist primär aus dunklem, gestanztem Stahlblech gefertigt, was ihr ein sehr funktionales, industrielles Aussehen verleiht. Sie verfügt über einen einfachen Pistolengriff, oft aus dunklem Polymer, und besitzt keinen traditionellen Holzschaft. Ihre markantesten Merkmale sind der gerade, vertikale Magazinschacht und die Metallschulterstütze, die nach unten und unter das Gehäuse klappt.
\end{pexample}

    {\small Feuerstoß}\\
    {\small Einzelschuss}\\

\begin{tabular}{@{}ll@{}}
    \textbf{\trans{weight_label}:} & 3.97 kg \\
    \textbf{\trans{price_label}:} & 1600 \\
    \textbf{Reichweite} & 40\\
    \textbf{Verborgenheit} & 4\\
    \textbf{Kapaziät} & 32\\
    \textbf{Schadenspotential} & 3\\
\end{tabular}

\wbif{phasesix}{
    \vspace{3mm}
\faIcon{flag}\hspace{0.5em}\faIcon{plane}\hspace{0.5em}\faIcon{microchip}\hspace{0.5em}}{}



\subsection*{Tommy Gun}

\begin{pexample}
Ein schweres, dunkles Stahlgehäuse ist mit einem polierten Holzschaft, einem hinteren Pistolengriff und einem markanten vertikalen Vordergriff verbunden. Ihr bekanntestes Merkmal ist das große, runde Trommelmagazin, das unter dem Verschluss eingerastet ist, obwohl sie auch gerade Stangenmagazine aufnimmt. Der relativ kurze Lauf ist oft von hervorstehenden Kühlrippen umschlossen.
\end{pexample}

    {\small Feuerstoß}\\

\begin{tabular}{@{}ll@{}}
    \textbf{\trans{weight_label}:} & 4.90 kg \\
    \textbf{\trans{price_label}:} & 2000 \\
    \textbf{Durchschlag} & 1\\
    \textbf{Reichweite} & 40\\
    \textbf{Verborgenheit} & 5\\
    \textbf{Kapaziät} & 50\\
    \textbf{Schadenspotential} & 3\\
\end{tabular}

\wbif{phasesix}{
    \vspace{3mm}
\faIcon{flag}\hspace{0.5em}\faIcon{plane}\hspace{0.5em}\faIcon{microchip}\hspace{0.5em}}{}



\subsection*{Uzi}

\begin{pexample}
Diese kompakte Feuerwaffe ist aus kantigem, gestanztem Schwarzmetall gefertigt, was ihr ein deutlich kastenförmiges Aussehen verleiht. Ihr markantestes Merkmal ist das lange, rechteckige Magazin, das direkt in den zentralen Pistolengriff eingeführt wird. Ein einfacher, klappbarer Metall-Schaft ist am hinteren Teil des Gehäuses angebracht, und der Lauf ist außergewöhnlich kurz.
\end{pexample}

    {\small Einzelschuss}\\
    {\small Feuerstoß}\\

\begin{tabular}{@{}ll@{}}
    \textbf{\trans{weight_label}:} & 3.60 kg \\
    \textbf{\trans{price_label}:} & 1300 \\
    \textbf{Reichweite} & 40\\
    \textbf{Verborgenheit} & 3\\
    \textbf{Kapaziät} & 32\\
    \textbf{Schadenspotential} & 2\\
\end{tabular}

\wbif{phasesix}{
    \vspace{3mm}
\faIcon{plane}\hspace{0.5em}\faIcon{microchip}\hspace{0.5em}}{}



\section{Schrotflinten}



\subsection*{Doppelläufige Schrotflinte}

\begin{pexample}
Zwei dicke, parallele Stahlläufe, die nebeneinander (oder übereinander) montiert sind, bilden den Kern dieser Feuerwaffe. Diese Laufbaugruppe ist mit einem massiven Holzschaft und Vorderschaft verbunden. Ein markanter Hebel oben auf dem Gehäuse oder eine seitliche Verriegelung dient als Freigabe, wodurch die Läufe am Verschluss nach unten geschwenkt werden können.
\end{pexample}

    {\small Einzelschuss}\\

\begin{tabular}{@{}ll@{}}
    \textbf{\trans{weight_label}:} & 2.00 kg \\
    \textbf{\trans{price_label}:} & 1200 \\
    \textbf{Durchschlag} & 1\\
    \textbf{Reichweite} & 15\\
    \textbf{Verborgenheit} & 6\\
    \textbf{Kapaziät} & 2\\
    \textbf{Schadenspotential} & 3\\
\end{tabular}

\wbif{phasesix}{
    \vspace{3mm}
\faIcon{flag}\hspace{0.5em}\faIcon{plane}\hspace{0.5em}\faIcon{microchip}\hspace{0.5em}}{}



\subsection*{Pump Action}

\begin{pexample}
Ein langes Röhrenmagazin verläuft parallel unter einem einzelnen, schweren Stahllauf. Diese Feuerwaffe ist sofort an ihrem verschiebbaren Vorderschaft (der „Pump“) aus Holz oder Polymer zu erkennen. Dieser Griff muss manuell vor- und zurückbewegt werden, um eine abgeschossene Hülse auszuwerfen und eine neue Patrone aus dem Magazin zuzuführen.
\end{pexample}

    {\small Einzelschuss}\\

\begin{tabular}{@{}ll@{}}
    \textbf{\trans{weight_label}:} & 2.00 kg \\
    \textbf{\trans{price_label}:} & 800 \\
    \textbf{Durchschlag} & 1\\
    \textbf{Reichweite} & 15\\
    \textbf{Verborgenheit} & 4\\
    \textbf{Kapaziät} & 6\\
    \textbf{Schadenspotential} & 2\\
\end{tabular}

\wbif{phasesix}{
    \vspace{3mm}
\faIcon{microchip}\hspace{0.5em}}{}



\subsection*{Abgesägte Schrotflinte}

\begin{pexample}
Zwei weite, parallele Stahlläufe sind mit einem einfachen Holzschaft und Griff verbunden. Die Waffe ist durch ihre außergewöhnlich kurze Länge gekennzeichnet. Sowohl die Metallläufe als auch der Holzschaft wurden grob abgesägt und weisen raue, unfertige Enden auf.
\end{pexample}

    {\small Einzelschuss}\\

\begin{tabular}{@{}ll@{}}
    \textbf{\trans{weight_label}:} & 1.30 kg \\
    \textbf{\trans{price_label}:} & 800 \\
    \textbf{Reichweite} & 15\\
    \textbf{Verborgenheit} & 3\\
    \textbf{Vorbereitung} & 0\\
    \textbf{Kapaziät} & 2\\
    \textbf{Schadenspotential} & 3\\
\end{tabular}

\wbif{phasesix}{
    \vspace{3mm}
\faIcon{flag}\hspace{0.5em}\faIcon{plane}\hspace{0.5em}\faIcon{microchip}\hspace{0.5em}}{}



\section{Wurfwaffen}



\subsection*{Wurfspeer}

\begin{pexample}
Ein schlanker Schaft aus flexiblem, leichtem Holz bildet den Körper dieses Speers. Sein Kopf ist eine kleine, scharfe Blattklinge oder ein simpler Dorn, dem der schwere, breite Kopf eines Militärspeers fehlt. Die Waffe ist merklich vor dem Griff ausbalanciert, eindeutig für den Wurf statt für den Nahkampf ausgelegt.
\end{pexample}

    {\small Werfen}\\

\begin{tabular}{@{}ll@{}}
    \textbf{\trans{weight_label}:} & 2.00 kg \\
    \textbf{\trans{price_label}:} & 300 \\
    \textbf{Durchschlag} & 1\\
    \textbf{Reichweite} & 20\\
    \textbf{Verborgenheit} & 6\\
    \textbf{Kapaziät} & 1\\
    \textbf{Schadenspotential} & 1\\
\end{tabular}

\wbif{phasesix}{
    \vspace{3mm}
\faIcon{eye}\hspace{0.5em}\faIcon{chess-rook}\hspace{0.5em}\faIcon{umbrella}\hspace{0.5em}}{}



\subsection*{Wurfaxt}

\begin{pexample}
Eine ausgeprägte Kopflastigkeit, eindeutig auf Rotation ausgelegt, definiert diese kleine Axt. Sie besitzt einen kurzen, glatten Holzschaft, der für einen einhändigen Griff ausgelegt ist. Der Metallkopf ist einschneidig und verfügt auf der Rückseite oft über einen geschärften Nacken oder Dorn, was sie von einem einfachen Handbeil unterscheidet.
\end{pexample}

    {\small Werfen}\\

\begin{tabular}{@{}ll@{}}
    \textbf{\trans{weight_label}:} & 1.00 kg \\
    \textbf{\trans{price_label}:} & 25 \\
    \textbf{Reichweite} & 10\\
    \textbf{Verborgenheit} & 3\\
    \textbf{Kapaziät} & 1\\
    \textbf{Schadenspotential} & 3\\
\end{tabular}

\wbif{phasesix}{
    \vspace{3mm}
\faIcon{cube}\hspace{0.5em}}{}



\subsection*{Shuriken}

\begin{pexample}
Der Shuriken ist ein kleines, dünnes, scheiben- oder sternförmiges Stück flachen, dunklen Metalls. Sie weist mehrere geschärfte Spitzen oder Kanten auf, die von einem zentralen Punkt oder Loch ausgehen.
\end{pexample}

    {\small Werfen}\\

\begin{tabular}{@{}ll@{}}
    \textbf{\trans{weight_label}:} & 0.20 kg \\
    \textbf{\trans{price_label}:} & 50 \\
    \textbf{Reichweite} & 10\\
    \textbf{Vorbereitung} & 0\\
    \textbf{Kapaziät} & 1\\
    \textbf{Schadenspotential} & 2\\
\end{tabular}

\wbif{phasesix}{
    \vspace{3mm}
\faIcon{umbrella}\hspace{0.5em}\faIcon{flag}\hspace{0.5em}\faIcon{plane}\hspace{0.5em}\faIcon{microchip}\hspace{0.5em}}{}



\section{Maschinengewehre}



\subsection*{Gatling}

\begin{pexample}
Ein Bündel aus langen, schweren Stahlläufen dominiert die Gatling, angeordnet in einem runden Gehäuse um eine Zentralachse. Diese Konstruktion ist auf einem schweren Gestell montiert und weist an der Seite eine markante, handbetriebene Kurbel auf. Ein großer Kasten oder eine Trommel, ausgelegt für die Munitionsaufnahme, ist typischerweise oben auf dem Verschlussgehäuse der Waffe angebracht.
\end{pexample}

    {\small Feuerstoß}\\

\begin{tabular}{@{}ll@{}}
    \textbf{\trans{weight_label}:} & 10.00 kg \\
    \textbf{\trans{price_label}:} & 8000 \\
    \textbf{Durchschlag} & 1\\
    \textbf{Reichweite} & 300\\
    \textbf{Verborgenheit} & 10\\
    \textbf{Vorbereitung} & 3\\
    \textbf{Kapaziät} & 100\\
    \textbf{Nachladeaktionen} & 3\\
    \textbf{Schadenspotential} & 4\\
\end{tabular}

\wbif{phasesix}{
    \vspace{3mm}
\faIcon{flag}\hspace{0.5em}}{}



\subsection*{M2 Browning}

\begin{pexample}
Diese enorme automatische Waffe basiert auf einem massiven, rechteckigen Stahlgehäuse. Sie verfügt über einen außergewöhnlich langen, schwerkalibrigen Lauf, der oft von einem dicken, perforierten Mantel umschlossen wird. Sie ist für die Zufuhr durch einen zerfallenden Metall-Munitionsgurt ausgelegt und wird durch ihr Paar Spatengriffe mit einem Schmetterlingsabzug am Heck charakterisiert.
\end{pexample}

    {\small Einzelschuss}\\
    {\small Feuerstoß}\\

\begin{tabular}{@{}ll@{}}
    \textbf{\trans{weight_label}:} & 34.20 kg \\
    \textbf{\trans{price_label}:} & 7000 \\
    \textbf{Durchschlag} & 1\\
    \textbf{Reichweite} & 100\\
    \textbf{Verborgenheit} & 10\\
    \textbf{Vorbereitung} & 2\\
    \textbf{Kapaziät} & 200\\
    \textbf{Schadenspotential} & 4\\
\end{tabular}

\wbif{phasesix}{
    \vspace{3mm}
\faIcon{flag}\hspace{0.5em}\faIcon{plane}\hspace{0.5em}\faIcon{microchip}\hspace{0.5em}}{}



\subsection*{M60}

\begin{pexample}
A bulky, stamped-steel receiver forms the core of this modern automatic weapon. It features a heavy barrel (often with an attached handle), a distinctive sloping shoulder stock, and a pistol grip. A folding bipod is attached near the front, and the weapon is designed to be fed by an ammunition belt from the left side.
\end{pexample}

    {\small Feuerstoß}\\
    {\small Einzelschuss}\\

\begin{tabular}{@{}ll@{}}
    \textbf{\trans{weight_label}:} & 10.50 kg \\
    \textbf{\trans{price_label}:} & 8900 \\
    \textbf{Durchschlag} & 1\\
    \textbf{Reichweite} & 80\\
    \textbf{Verborgenheit} & 9\\
    \textbf{Vorbereitung} & 3\\
    \textbf{Kapaziät} & 250\\
    \textbf{Schadenspotential} & 3\\
\end{tabular}

\wbif{phasesix}{
    \vspace{3mm}
\faIcon{plane}\hspace{0.5em}\faIcon{microchip}\hspace{0.5em}}{}



\subsection*{MG34}

\begin{pexample}
Ein langer, schlanker Lauf, der von einem markanten, runden und mit vielen Löchern perforierten Mantel umgeben ist, definiert diese automatische Waffe. Sie ist mit einem dunklen, gefrästen Stahlgehäuse, einem polierten Holzschaft und einem einfachen Pistolengriff gebaut. Ein klappbares Zweibein ist oft am Laufmantel festgeklemmt, der für einen schnellen Wechsel ausgelegt ist.
\end{pexample}

    {\small Feuerstoß}\\
    {\small Einzelschuss}\\

\begin{tabular}{@{}ll@{}}
    \textbf{\trans{weight_label}:} & 12.10 kg \\
    \textbf{\trans{price_label}:} & 5500 \\
    \textbf{Durchschlag} & 2\\
    \textbf{Reichweite} & 80\\
    \textbf{Verborgenheit} & 9\\
    \textbf{Vorbereitung} & 3\\
    \textbf{Kapaziät} & 200\\
    \textbf{Schadenspotential} & 4\\
\end{tabular}

\wbif{phasesix}{
    \vspace{3mm}
\faIcon{flag}\hspace{0.5em}\faIcon{plane}\hspace{0.5em}\faIcon{microchip}\hspace{0.5em}}{}



\subsection*{MG42}

\begin{pexample}
Diese moderne automatische Waffe basiert auf einem langen Gehäuse aus Stahlblech und verfügt über einen markanten, perforierten Laufmantel zur Hitzeregulierung. Ein klappbares Zweibein ist nahe der Mündung angebracht, und sie ist mit einem einfachen Pistolengriff sowie einem flossenförmigen Holz- oder Polymerschaft ausgestattet. Ein markanter Zuführdeckel an der Seite des Gehäuses weist auf den Gurtzuführungs-Mechanismus hin.
\end{pexample}

    {\small Feuerstoß}\\

\begin{tabular}{@{}ll@{}}
    \textbf{\trans{weight_label}:} & 10.60 kg \\
    \textbf{\trans{price_label}:} & 6000 \\
    \textbf{Durchschlag} & 2\\
    \textbf{Reichweite} & 80\\
    \textbf{Vorbereitung} & 2\\
    \textbf{Kapaziät} & 200\\
    \textbf{Schadenspotential} & 4\\
    \textbf{Verborgenheit} & 9\\
\end{tabular}

\wbif{phasesix}{
    \vspace{3mm}
\faIcon{flag}\hspace{0.5em}\faIcon{plane}\hspace{0.5em}\faIcon{microchip}\hspace{0.5em}}{}



\subsection*{PKM}

\begin{pexample}
Diese moderne automatische Waffe ist auf einem langen Gehäuse aus Stahlblech aufgebaut und mit einem schweren, gerippten Lauf für Dauerfeuer ausgestattet. Sie verfügt über einen markanten, durchbrochenen Schaft aus Holz oder Polymer, einen Pistolengriff und ein klappbares Zweibein nahe der Mündung. Die Munition wird über einen zerfallenden Metallgurt zugeführt, der sich oft in einem Kasten befindet, der direkt unter dem Gehäuse eingehängt ist.
\end{pexample}

    {\small Feuerstoß}\\

\begin{tabular}{@{}ll@{}}
    \textbf{\trans{weight_label}:} & 7.50 kg \\
    \textbf{\trans{price_label}:} & 8000 \\
    \textbf{Durchschlag} & 1\\
    \textbf{Reichweite} & 80\\
    \textbf{Verborgenheit} & 9\\
    \textbf{Kapaziät} & 200\\
    \textbf{Schadenspotential} & 3\\
\end{tabular}

\wbif{phasesix}{
    \vspace{3mm}
\faIcon{flag}\hspace{0.5em}\faIcon{plane}\hspace{0.5em}\faIcon{microchip}\hspace{0.5em}}{}



\section{Schwere Waffen}



\subsection*{Flammenwerfer}

\begin{pexample}
Diese Waffe besteht aus einem oder zwei großen Metalltanks, die mit einem Gurtzeug auf dem Rücken getragen und durch einen dicken, verstärkten Schlauch mit einem langen, metallenen Strahlrohr verbunden sind. Das Strahlrohr verfügt über eine Abzugseinheit und eine Düse, aus der es einen Strahl entzündeten Treibstoffs aussendet. Der Flammenwerfer wirkt in einem 45-Grad-Kegel. Treffer werden vom Schützen unter allen Gegnern im Kegel verteilt und verursacht bei allen Getroffenen den Status Brennend 1.
\end{pexample}

    {\small Feuerstoß}\\

\begin{tabular}{@{}ll@{}}
    \textbf{\trans{weight_label}:} & 6.00 kg \\
    \textbf{\trans{price_label}:} & 10000 \\
    \textbf{Reichweite} & 10\\
    \textbf{Verborgenheit} & 7\\
    \textbf{Nachladeaktionen} & 20\\
    \textbf{Schadenspotential} & 8\\
    \textbf{Vorbereitung} & 3\\
    \textbf{Kapaziät} & 30\\
    \textbf{Brennend} & 1\\
    \textbf{Kegel} & 45\\
\end{tabular}

\wbif{phasesix}{
    \vspace{3mm}
\faIcon{flag}\hspace{0.5em}\faIcon{plane}\hspace{0.5em}\faIcon{microchip}\hspace{0.5em}\faIcon{fire}\hspace{0.5em}}{}



\subsection*{HK GMW}

\begin{pexample}
Diese schwere, kastenförmige Automatikwaffe ist ein Maschinengranatwerfer im Kaliber 40x53mm, der auf einem soliden Metallgehäuse aufgebaut ist. Er verfügt über eine markante, seitliche Gurtzuführung für 32 Schuss Munition, was ihn von Einzelschuss-Granatwerfern unterscheidet. Die gesamte Baugruppe ist für Dauerfeuer ausgelegt und wird oft auf einem Dreibein oder Fahrzeug montiert.
\end{pexample}

    {\small Einzelschuss}\\

\begin{tabular}{@{}ll@{}}
    \textbf{\trans{weight_label}:} & 29.00 kg \\
    \textbf{\trans{price_label}:} & 2000 \\
    \textbf{Durchschlag} & 1\\
    \textbf{Reichweite} & 40\\
    \textbf{Verborgenheit} & 5\\
    \textbf{Vorbereitung} & 2\\
    \textbf{Kapaziät} & 32\\
    \textbf{Nachladeaktionen} & 3\\
    \textbf{Schadenspotential} & 4\\
\end{tabular}

\wbif{phasesix}{
    \vspace{3mm}
\faIcon{microchip}\hspace{0.5em}}{}



\subsection*{M79 (Granatwerfer)}

\begin{pexample}
The M79 grenade launcher, also known as a thumper, thump gun, bloop tube or blooper because of its special muzzle sound, was introduced into the US Army in 1961. It is a shoulder-fired system for firing 40 mm grenades and is designed to cover the gap between 50 m (maximum throw distance of hand grenades) and 300 m (minimum mortar distance).
\end{pexample}

    {\small Einzelschuss}\\

\begin{tabular}{@{}ll@{}}
    \textbf{\trans{weight_label}:} & 3.00 kg \\
    \textbf{\trans{price_label}:} & 1200 \\
    \textbf{Durchschlag} & 1\\
    \textbf{Reichweite} & 150\\
    \textbf{Verborgenheit} & 4\\
    \textbf{Kapaziät} & 1\\
    \textbf{Schadenspotential} & 4\\
\end{tabular}

\wbif{phasesix}{
    \vspace{3mm}
\faIcon{plane}\hspace{0.5em}\faIcon{microchip}\hspace{0.5em}}{}



\end{multicols}